\documentclass[12pt,aspectratio=169]{beamer}
\usepackage[T1]{fontenc}
\usepackage{lmodern}
\usepackage{amsmath,amssymb,array,booktabs,tabularx,tikz}
% Add language support appropriate to the source, if needed.
% Keep this preamble inline in every delivered deck.
\definecolor{GEInk}{HTML}{132A3E}
\definecolor{GETeal}{HTML}{17675E}
\definecolor{GEFooter}{HTML}{EAF2F1}
\setbeamercolor{background canvas}{bg=white}
\setbeamercolor{normal text}{fg=GEInk,bg=white}
\setbeamercolor{structure}{fg=GETeal}
\setbeamercolor{frametitle}{fg=GEInk}
\setbeamercolor{GE footer}{fg=GETeal,bg=GEFooter}
\setbeamerfont{normal text}{family=\sffamily,size=\normalsize}
\setbeamerfont{title}{size=\huge,series=\bfseries}
\setbeamerfont{subtitle}{size=\normalsize,series=\bfseries}
\setbeamerfont{frametitle}{size=\Large,series=\bfseries}
\setbeamerfont{framesubtitle}{size=\normalsize,series=\mdseries}
\setbeamerfont{GE footer}{size=\footnotesize}
\setbeamersize{text margin left=9mm,text margin right=9mm}
\setbeamertemplate{navigation symbols}{}
\setbeamertemplate{headline}{}
\setbeamercovered{invisible}
\setbeamertemplate{frametitle}{%
  \vskip.5em
  {\usebeamerfont{frametitle}\strut\insertframetitle\strut\par}%
  \ifx\insertframesubtitle\empty\else
    {\usebeamerfont{framesubtitle}\strut\insertframesubtitle\strut\par}%
  \fi
  \vskip.25em
}
\setbeamertemplate{footline}{%
  \begin{beamercolorbox}[wd=\paperwidth,ht=3ex,dp=1.3ex,
    leftskip=9mm,rightskip=9mm]{GE footer}%
    \usebeamerfont{GE footer}\insertsectionhead\hfill
    \insertframenumber/\inserttotalframenumber
  \end{beamercolorbox}%
}
\setlength{\parindent}{0pt}
\setlength{\parskip}{.35em}
\setlength{\tabcolsep}{4pt}
\renewcommand{\arraystretch}{1.22}
\AtBeginEnvironment{frame}{\raggedright}
\addtobeamertemplate{frametitle}{}{%
  \setlength{\abovedisplayskip}{6pt plus 2pt minus 2pt}%
  \setlength{\belowdisplayskip}{6pt plus 2pt minus 2pt}%
  \setlength{\abovedisplayshortskip}{6pt plus 2pt minus 2pt}%
  \setlength{\belowdisplayshortskip}{6pt plus 2pt minus 2pt}%
}






\usetikzlibrary{shapes.geometric,decorations.pathreplacing,arrows.meta}
\usepackage{cancel}
\definecolor{CircleBlue}{HTML}{2879B9}
\definecolor{TriangleOrange}{HTML}{DD8428}
\definecolor{SquareGreen}{HTML}{398553}
\definecolor{DiamondPurple}{HTML}{8954A2}
% Source map: permutations1.tex, Introduction: Counting Choices ONLY.
% Stop before Worked Examples. The source practice is preserved unchanged.
% Case 1: multiplication paragraph; case 2: factorial bullet;
% case 3: ordered positions bullet; case 4: unordered groups bullet;
% case 5: indistinguishable copies bullet; closing frames synthesize these.
% The fifteen concrete illustrations and opening are specified by the user.
% Scope adaptation: no official criteria, invented metadata or highlight registry.
% Colors identify supplied shape types, not mathematical highlight roles.
% Shapes, identifiers, outlines and text preserve meaning in grayscale.
\tikzset{every picture/.style={draw=GEInk,text=GEInk,line width=.7pt},
 group/.style={draw=GEInk,rounded corners=2mm,line width=1pt},
 object/.style={draw=GEInk,minimum size=3.4mm,inner sep=0pt},
 C/.style={object,circle,fill=CircleBlue},
 T/.style={object,regular polygon,regular polygon sides=3,fill=TriangleOrange},
 S/.style={object,rectangle,fill=SquareGreen},
 D/.style={object,diamond,aspect=1,fill=DiamondPurple}}
% Object identifiers are ABOVE objects. Position/job labels are BELOW rows.
\newcommand{\Obj}[4]{\node[#1] at (#2,#3) {};\node[font=\small] at (#2,{#3+.30}) {#4};}
\newcommand{\Q}[1]{\textbf{#1}\par}
\newcommand{\DiagramStart}{\par\centering\begin{tikzpicture}
 \path[use as bounding box] (-.05,-.55) rectangle (13.55,4.05);}
\newcommand{\DiagramEnd}{\end{tikzpicture}\par\raggedright}
\begin{document}
\section{Five cases}
% Frame 1. Teaching note: Opening overview deliberately has no formulas. Return to the question for every example.
\begin{frame}[t]{Counting choices: five cases}
\Q{What counts as one different outcome?}
\DiagramStart
\node[anchor=west] at (0.1,3.45) {1. Choose one from each input group};\Obj{C}{10.5}{3.3800000000000003}{1}\Obj{T}{11.04}{3.3800000000000003}{1}\draw (10.25,3.12) rectangle (10.75,3.85);\draw (10.79,3.12) rectangle (11.29,3.85);\node[anchor=west] at (0.1,2.66) {2. Arrange all distinct objects};\Obj{C}{10.5}{2.5900000000000003}{1}\Obj{C}{11.04}{2.5900000000000003}{2}\Obj{C}{11.58}{2.5900000000000003}{3}\node[anchor=west] at (0.1,1.87) {3. Fill some ordered positions};\Obj{C}{10.5}{1.8}{1}\Obj{C}{11.04}{1.8}{2}\node[font=\scriptsize] at (10.5,1.53) {Reader};\node[font=\scriptsize] at (11.04,1.53) {Recorder};\node[anchor=west] at (0.1,1.08) {4. Choose a group without roles};\Obj{C}{10.5}{1.01}{1}\Obj{C}{11.04}{1.01}{2}\draw (10.77,1.03) ellipse (.65 and .43);\node[anchor=west] at (0.1,0.29000000000000004) {5. Arrange indistinguishable copies};\Obj{C}{10.5}{0.22000000000000003}{}\Obj{C}{11.04}{0.22000000000000003}{}\Obj{T}{11.58}{0.22000000000000003}{}
\DiagramEnd
\begin{overlayarea}{\linewidth}{1.15cm}\end{overlayarea}
\end{frame}

% Frame 2. Teaching note: Misconception: multiplying object labels. Count choices, not the identifying numbers.
\begin{frame}[t]{How to read our objects}
\Q{Does object 2 mean a value of two?}
\DiagramStart
\Obj{C}{4.8}{2.6}{1}\Obj{C}{5.9}{2.6}{2}\Obj{C}{7.0}{2.6}{3}\node[] at (6,1.4) {Same shape and color; different identities.};
\DiagramEnd
\begin{overlayarea}{\linewidth}{1.15cm}\only<2>{No. A number above a shape identifies that object.\par}
\only<3->{Position labels go below the row. Shapes also work without color.\par}\end{overlayarea}
\end{frame}

% Frame 3. Teaching note: Collect predictions before revealing outcomes. Numbers identify objects; never multiply the identifiers.
\begin{frame}[t]{1A. One circle, then one triangle}
\Q{Predict: how many different outcomes are possible?}
\DiagramStart
\Obj{C}{2.0}{2.7}{1}\Obj{C}{2.75}{2.7}{2}\node[] at (2.4,1.8) {Circle choices};\Obj{T}{7.0}{2.7}{1}\Obj{T}{7.75}{2.7}{2}\Obj{T}{8.5}{2.7}{3}\node[] at (7.4,1.8) {Triangle choices};\node[] at (6.75,0.5) {One outcome: one circle followed by one triangle.};
\DiagramEnd
\begin{overlayarea}{\linewidth}{1.15cm}\end{overlayarea}
\end{frame}

\section{1. Choose from input groups}
% Frame 4. Teaching note: Reveal rows in order. At overlay 6 replace row boundaries with column boundaries; do not move outcomes.
\begin{frame}[t]{1A. The same outcomes, two organizations}
\Q{How many outcomes are in each visual group?}
\DiagramStart
\visible<2->{\Obj{C}{1.5}{2.8}{1}\Obj{T}{2.05}{2.8}{1}}\visible<2->{\Obj{C}{6.75}{2.8}{1}\Obj{T}{7.3}{2.8}{2}}\visible<2->{\Obj{C}{12.0}{2.8}{1}\Obj{T}{12.55}{2.8}{3}}\visible<3->{\Obj{C}{1.5}{0.3999999999999999}{2}\Obj{T}{2.05}{0.3999999999999999}{1}}\visible<3->{\Obj{C}{6.75}{0.3999999999999999}{2}\Obj{T}{7.3}{0.3999999999999999}{2}}\visible<3->{\Obj{C}{12.0}{0.3999999999999999}{2}\Obj{T}{12.55}{0.3999999999999999}{3}}\visible<4-5>{\draw[group] (1.05,2.3499999999999996) rectangle (12.95,3.3699999999999997);}\visible<4-5>{\draw[group] (1.05,-0.0500000000000001) rectangle (12.95,0.9699999999999999);}\visible<6->{\draw[group] (1.05,-0.05) rectangle (2.5,3.37);}\visible<6->{\draw[group] (6.3,-0.05) rectangle (7.75,3.37);}\visible<6->{\draw[group] (11.55,-0.05) rectangle (13.0,3.37);}
\DiagramEnd
\begin{overlayarea}{\linewidth}{1.15cm}\only<4-5>{Two row groups; three outcomes in each.\par}
\only<5>{$2\times3=6$ outcomes.\par}
\only<6>{Three column groups; two outcomes in each.\par}
\only<7->{$3\times2=6$. These are the same six outcomes.\par}\end{overlayarea}
\end{frame}

% Frame 5. Teaching note: Collect predictions before revealing outcomes. Numbers identify objects; never multiply the identifiers.
\begin{frame}[t]{1B. One circle, one triangle, one square}
\Q{Predict: how many different outcomes are possible?}
\DiagramStart
\Obj{C}{1.0}{2.7}{1}\node[] at (1.4,1.8) {Circle};\Obj{T}{4.0}{2.7}{1}\Obj{T}{4.75}{2.7}{2}\Obj{T}{5.5}{2.7}{3}\node[] at (4.4,1.8) {Triangles};\Obj{S}{9.0}{2.7}{1}\Obj{S}{9.75}{2.7}{2}\Obj{S}{10.5}{2.7}{3}\Obj{S}{11.25}{2.7}{4}\node[] at (9.4,1.8) {Squares};\node[] at (6.75,0.5) {One outcome: circle, triangle, square, in that order.};
\DiagramEnd
\begin{overlayarea}{\linewidth}{1.15cm}\end{overlayarea}
\end{frame}

% Frame 6. Teaching note: Only one grouping view at a time. The factor 1 never creates another branch.
\begin{frame}[t]{1B. All outcomes keep the same circle}
\Q{Which choices change across rows and columns?}
\DiagramStart
\visible<2->{\Obj{C}{1.5}{2.8}{1}\Obj{T}{2.05}{2.8}{1}\Obj{S}{2.6}{2.8}{1}}\visible<2->{\Obj{C}{5.0}{2.8}{1}\Obj{T}{5.55}{2.8}{1}\Obj{S}{6.1}{2.8}{2}}\visible<2->{\Obj{C}{8.5}{2.8}{1}\Obj{T}{9.05}{2.8}{1}\Obj{S}{9.6}{2.8}{3}}\visible<2->{\Obj{C}{12.0}{2.8}{1}\Obj{T}{12.55}{2.8}{1}\Obj{S}{13.1}{2.8}{4}}\visible<3->{\Obj{C}{1.5}{1.5999999999999999}{1}\Obj{T}{2.05}{1.5999999999999999}{2}\Obj{S}{2.6}{1.5999999999999999}{1}}\visible<3->{\Obj{C}{5.0}{1.5999999999999999}{1}\Obj{T}{5.55}{1.5999999999999999}{2}\Obj{S}{6.1}{1.5999999999999999}{2}}\visible<3->{\Obj{C}{8.5}{1.5999999999999999}{1}\Obj{T}{9.05}{1.5999999999999999}{2}\Obj{S}{9.6}{1.5999999999999999}{3}}\visible<3->{\Obj{C}{12.0}{1.5999999999999999}{1}\Obj{T}{12.55}{1.5999999999999999}{2}\Obj{S}{13.1}{1.5999999999999999}{4}}\visible<4->{\Obj{C}{1.5}{0.3999999999999999}{1}\Obj{T}{2.05}{0.3999999999999999}{3}\Obj{S}{2.6}{0.3999999999999999}{1}}\visible<4->{\Obj{C}{5.0}{0.3999999999999999}{1}\Obj{T}{5.55}{0.3999999999999999}{3}\Obj{S}{6.1}{0.3999999999999999}{2}}\visible<4->{\Obj{C}{8.5}{0.3999999999999999}{1}\Obj{T}{9.05}{0.3999999999999999}{3}\Obj{S}{9.6}{0.3999999999999999}{3}}\visible<4->{\Obj{C}{12.0}{0.3999999999999999}{1}\Obj{T}{12.55}{0.3999999999999999}{3}\Obj{S}{13.1}{0.3999999999999999}{4}}\visible<5-6>{\draw[group] (1.05,-0.06) rectangle (13.48,3.38);}\visible<7-8>{\draw[group] (1.05,2.3699999999999997) rectangle (13.48,3.3499999999999996);}\visible<7-8>{\draw[group] (1.05,1.17) rectangle (13.48,2.15);}\visible<7-8>{\draw[group] (1.05,-0.030000000000000082) rectangle (13.48,0.95);}\visible<9-10>{\draw[group] (1.05,-0.06) rectangle (2.95,3.38);}\visible<9-10>{\draw[group] (4.55,-0.06) rectangle (6.45,3.38);}\visible<9-10>{\draw[group] (8.05,-0.06) rectangle (9.95,3.38);}\visible<9-10>{\draw[group] (11.55,-0.06) rectangle (13.45,3.38);}
\DiagramEnd
\begin{overlayarea}{\linewidth}{1.15cm}\only<5-6>{One outer group: there is only one circle choice.\par}
\only<7-8>{Three row groups: one for each triangle; four outcomes each.\par}
\only<9-10>{Four column groups: one for each square; three outcomes each.\par}
\only<11->{$1\times3\times4=12$. The circle step offers no alternative.\par}\end{overlayarea}
\end{frame}

% Frame 7. Teaching note: Collect predictions before revealing outcomes. Numbers identify objects; never multiply the identifiers.
\begin{frame}[t]{1C. Four input groups}
\Q{Predict: how many different outcomes are possible?}
\DiagramStart
\Obj{C}{0.7}{2.7}{1}\Obj{C}{1.45}{2.7}{2}\node[] at (1.1,1.8) {Circles};\Obj{T}{3.7}{2.7}{1}\Obj{T}{4.45}{2.7}{2}\node[] at (4.1000000000000005,1.8) {Triangles};\Obj{S}{7.0}{2.7}{1}\Obj{S}{7.75}{2.7}{2}\node[] at (7.4,1.8) {Squares};\Obj{D}{10.5}{2.7}{1}\Obj{D}{11.25}{2.7}{2}\node[] at (10.9,1.8) {Diamonds};\node[] at (6.75,0.5) {One outcome: circle, triangle, square, diamond.};
\DiagramEnd
\begin{overlayarea}{\linewidth}{1.15cm}\end{overlayarea}
\end{frame}

% Frame 8. Teaching note: Panel rows fix circle; panel columns fix triangle. Inside panels rows vary square and columns vary diamond.
\begin{frame}[t]{1C. Four panels, each with its own fixed start}
\Q{How many outcomes belong to each panel?}
\DiagramStart
\node[anchor=west,font=\small] at (0.6,3.43) {Fixed:};\Obj{C}{1.9}{3.43}{1}\Obj{T}{2.44}{3.43}{1}\visible<2->{\Obj{C}{0.6}{2.7}{1}\Obj{T}{1.06}{2.7}{1}\Obj{S}{1.52}{2.7}{1}\Obj{D}{1.98}{2.7}{1}}\visible<2->{\Obj{C}{3.4}{2.7}{1}\Obj{T}{3.86}{2.7}{1}\Obj{S}{4.32}{2.7}{1}\Obj{D}{4.78}{2.7}{2}}\visible<2->{\Obj{C}{0.6}{1.8800000000000003}{1}\Obj{T}{1.06}{1.8800000000000003}{1}\Obj{S}{1.52}{1.8800000000000003}{2}\Obj{D}{1.98}{1.8800000000000003}{1}}\visible<2->{\Obj{C}{3.4}{1.8800000000000003}{1}\Obj{T}{3.86}{1.8800000000000003}{1}\Obj{S}{4.32}{1.8800000000000003}{2}\Obj{D}{4.78}{1.8800000000000003}{2}}\visible<6-8>{\draw[group] (0.3,1.52) rectangle (5.3,3.9699999999999998);}\node[anchor=west,font=\small] at (7.3,3.43) {Fixed:};\Obj{C}{8.6}{3.43}{1}\Obj{T}{9.14}{3.43}{2}\visible<3->{\Obj{C}{7.3}{2.7}{1}\Obj{T}{7.76}{2.7}{2}\Obj{S}{8.22}{2.7}{1}\Obj{D}{8.68}{2.7}{1}}\visible<3->{\Obj{C}{10.1}{2.7}{1}\Obj{T}{10.56}{2.7}{2}\Obj{S}{11.02}{2.7}{1}\Obj{D}{11.48}{2.7}{2}}\visible<3->{\Obj{C}{7.3}{1.8800000000000003}{1}\Obj{T}{7.76}{1.8800000000000003}{2}\Obj{S}{8.22}{1.8800000000000003}{2}\Obj{D}{8.68}{1.8800000000000003}{1}}\visible<3->{\Obj{C}{10.1}{1.8800000000000003}{1}\Obj{T}{10.56}{1.8800000000000003}{2}\Obj{S}{11.02}{1.8800000000000003}{2}\Obj{D}{11.48}{1.8800000000000003}{2}}\visible<6-8>{\draw[group] (7.0,1.52) rectangle (12.0,3.9699999999999998);}\node[anchor=west,font=\small] at (0.6,1.43) {Fixed:};\Obj{C}{1.9}{1.43}{2}\Obj{T}{2.44}{1.43}{1}\visible<4->{\Obj{C}{0.6}{0.7}{2}\Obj{T}{1.06}{0.7}{1}\Obj{S}{1.52}{0.7}{1}\Obj{D}{1.98}{0.7}{1}}\visible<4->{\Obj{C}{3.4}{0.7}{2}\Obj{T}{3.86}{0.7}{1}\Obj{S}{4.32}{0.7}{1}\Obj{D}{4.78}{0.7}{2}}\visible<4->{\Obj{C}{0.6}{-0.12}{2}\Obj{T}{1.06}{-0.12}{1}\Obj{S}{1.52}{-0.12}{2}\Obj{D}{1.98}{-0.12}{1}}\visible<4->{\Obj{C}{3.4}{-0.12}{2}\Obj{T}{3.86}{-0.12}{1}\Obj{S}{4.32}{-0.12}{2}\Obj{D}{4.78}{-0.12}{2}}\visible<6-8>{\draw[group] (0.3,-0.48) rectangle (5.3,1.97);}\node[anchor=west,font=\small] at (7.3,1.43) {Fixed:};\Obj{C}{8.6}{1.43}{2}\Obj{T}{9.14}{1.43}{2}\visible<5->{\Obj{C}{7.3}{0.7}{2}\Obj{T}{7.76}{0.7}{2}\Obj{S}{8.22}{0.7}{1}\Obj{D}{8.68}{0.7}{1}}\visible<5->{\Obj{C}{10.1}{0.7}{2}\Obj{T}{10.56}{0.7}{2}\Obj{S}{11.02}{0.7}{1}\Obj{D}{11.48}{0.7}{2}}\visible<5->{\Obj{C}{7.3}{-0.12}{2}\Obj{T}{7.76}{-0.12}{2}\Obj{S}{8.22}{-0.12}{2}\Obj{D}{8.68}{-0.12}{1}}\visible<5->{\Obj{C}{10.1}{-0.12}{2}\Obj{T}{10.56}{-0.12}{2}\Obj{S}{11.02}{-0.12}{2}\Obj{D}{11.48}{-0.12}{2}}\visible<6-8>{\draw[group] (7.0,-0.48) rectangle (12.0,1.97);}
\DiagramEnd
\begin{overlayarea}{\linewidth}{1.15cm}\only<6>{Four panels. Inside each: two square choices and two diamond choices.\par}
\only<7>{Four outcomes per panel: $4\times4=16$.\par}
\only<8->{Each $4$ is $2\times2$, so $2\times2\times2\times2=16$.\par}\end{overlayarea}
\end{frame}

% Frame 9. Teaching note: Last-choice groups are nonadjacent cells. Use different boundary overlays, never simultaneous panel/row/cell outlines.
\begin{frame}[t]{1C. Regroup without moving any outcome}
\Q{Can the same sixteen outcomes make two groups of eight?}
\DiagramStart
\node[anchor=west,font=\small] at (0.6,3.43) {Fixed:};\Obj{C}{1.9}{3.43}{1}\Obj{T}{2.44}{3.43}{1}\visible<2->{\Obj{C}{0.6}{2.7}{1}\Obj{T}{1.06}{2.7}{1}\Obj{S}{1.52}{2.7}{1}\Obj{D}{1.98}{2.7}{1}}\visible<2->{\Obj{C}{3.4}{2.7}{1}\Obj{T}{3.86}{2.7}{1}\Obj{S}{4.32}{2.7}{1}\Obj{D}{4.78}{2.7}{2}}\visible<2->{\Obj{C}{0.6}{1.8800000000000003}{1}\Obj{T}{1.06}{1.8800000000000003}{1}\Obj{S}{1.52}{1.8800000000000003}{2}\Obj{D}{1.98}{1.8800000000000003}{1}}\visible<2->{\Obj{C}{3.4}{1.8800000000000003}{1}\Obj{T}{3.86}{1.8800000000000003}{1}\Obj{S}{4.32}{1.8800000000000003}{2}\Obj{D}{4.78}{1.8800000000000003}{2}}\node[anchor=west,font=\small] at (7.3,3.43) {Fixed:};\Obj{C}{8.6}{3.43}{1}\Obj{T}{9.14}{3.43}{2}\visible<3->{\Obj{C}{7.3}{2.7}{1}\Obj{T}{7.76}{2.7}{2}\Obj{S}{8.22}{2.7}{1}\Obj{D}{8.68}{2.7}{1}}\visible<3->{\Obj{C}{10.1}{2.7}{1}\Obj{T}{10.56}{2.7}{2}\Obj{S}{11.02}{2.7}{1}\Obj{D}{11.48}{2.7}{2}}\visible<3->{\Obj{C}{7.3}{1.8800000000000003}{1}\Obj{T}{7.76}{1.8800000000000003}{2}\Obj{S}{8.22}{1.8800000000000003}{2}\Obj{D}{8.68}{1.8800000000000003}{1}}\visible<3->{\Obj{C}{10.1}{1.8800000000000003}{1}\Obj{T}{10.56}{1.8800000000000003}{2}\Obj{S}{11.02}{1.8800000000000003}{2}\Obj{D}{11.48}{1.8800000000000003}{2}}\node[anchor=west,font=\small] at (0.6,1.43) {Fixed:};\Obj{C}{1.9}{1.43}{2}\Obj{T}{2.44}{1.43}{1}\visible<4->{\Obj{C}{0.6}{0.7}{2}\Obj{T}{1.06}{0.7}{1}\Obj{S}{1.52}{0.7}{1}\Obj{D}{1.98}{0.7}{1}}\visible<4->{\Obj{C}{3.4}{0.7}{2}\Obj{T}{3.86}{0.7}{1}\Obj{S}{4.32}{0.7}{1}\Obj{D}{4.78}{0.7}{2}}\visible<4->{\Obj{C}{0.6}{-0.12}{2}\Obj{T}{1.06}{-0.12}{1}\Obj{S}{1.52}{-0.12}{2}\Obj{D}{1.98}{-0.12}{1}}\visible<4->{\Obj{C}{3.4}{-0.12}{2}\Obj{T}{3.86}{-0.12}{1}\Obj{S}{4.32}{-0.12}{2}\Obj{D}{4.78}{-0.12}{2}}\node[anchor=west,font=\small] at (7.3,1.43) {Fixed:};\Obj{C}{8.6}{1.43}{2}\Obj{T}{9.14}{1.43}{2}\visible<5->{\Obj{C}{7.3}{0.7}{2}\Obj{T}{7.76}{0.7}{2}\Obj{S}{8.22}{0.7}{1}\Obj{D}{8.68}{0.7}{1}}\visible<5->{\Obj{C}{10.1}{0.7}{2}\Obj{T}{10.56}{0.7}{2}\Obj{S}{11.02}{0.7}{1}\Obj{D}{11.48}{0.7}{2}}\visible<5->{\Obj{C}{7.3}{-0.12}{2}\Obj{T}{7.76}{-0.12}{2}\Obj{S}{8.22}{-0.12}{2}\Obj{D}{8.68}{-0.12}{1}}\visible<5->{\Obj{C}{10.1}{-0.12}{2}\Obj{T}{10.56}{-0.12}{2}\Obj{S}{11.02}{-0.12}{2}\Obj{D}{11.48}{-0.12}{2}}\visible<6-7>{\draw[group] (0.3,1.52) rectangle (12,3.9699999999999998);}\visible<6-7>{\draw[group] (0.3,-0.48) rectangle (12,1.97);}\visible<8-9>{\draw[group,dashed] (0.35,2.45) rectangle (2.25,3.2);}\visible<10->{\draw[group] (3.15,2.45) rectangle (5.05,3.2);}\visible<8-9>{\draw[group,dashed] (0.35,1.6300000000000003) rectangle (2.25,2.3800000000000003);}\visible<10->{\draw[group] (3.15,1.6300000000000003) rectangle (5.05,2.3800000000000003);}\visible<8-9>{\draw[group,dashed] (7.05,2.45) rectangle (8.95,3.2);}\visible<10->{\draw[group] (9.85,2.45) rectangle (11.75,3.2);}\visible<8-9>{\draw[group,dashed] (7.05,1.6300000000000003) rectangle (8.95,2.3800000000000003);}\visible<10->{\draw[group] (9.85,1.6300000000000003) rectangle (11.75,2.3800000000000003);}\visible<8-9>{\draw[group,dashed] (0.35,0.44999999999999996) rectangle (2.25,1.2);}\visible<10->{\draw[group] (3.15,0.44999999999999996) rectangle (5.05,1.2);}\visible<8-9>{\draw[group,dashed] (0.35,-0.37) rectangle (2.25,0.38);}\visible<10->{\draw[group] (3.15,-0.37) rectangle (5.05,0.38);}\visible<8-9>{\draw[group,dashed] (7.05,0.44999999999999996) rectangle (8.95,1.2);}\visible<10->{\draw[group] (9.85,0.44999999999999996) rectangle (11.75,1.2);}\visible<8-9>{\draw[group,dashed] (7.05,-0.37) rectangle (8.95,0.38);}\visible<10->{\draw[group] (9.85,-0.37) rectangle (11.75,0.38);}
\DiagramEnd
\begin{overlayarea}{\linewidth}{1.15cm}\only<6-7>{Circle 1 gives eight; circle 2 gives eight. $2\times8=16$.\par}
\only<8-9>{The eight dashed boxes all end in diamond 1: one group of eight.\par}
\only<10->{The other eight all end in diamond 2: another group. $2\times8=16$.\par}\end{overlayarea}
\end{frame}

% Frame 10. Teaching note: Equal branch counts at each step are the condition; the factors need not equal one another.
\begin{frame}[t]{The multiplication principle}
\Q{When can we multiply choice counts?}
\DiagramStart
\node[] at (2,2.5) {Step 1};\node[] at (6.5,2.5) {Step 2};\node[] at (11,2.5) {Step 3};\node[] at (2,1.4) {2 choices};\node[] at (6.5,1.4) {3 choices each};\node[] at (11,1.4) {4 choices each};
\DiagramEnd
\begin{overlayarea}{\linewidth}{1.15cm}\only<2>{Every earlier choice must leave the same next-step choice count.\par}
\only<3>{The numbers can differ between steps: $2\times3\times4$.\par}
\only<4->{Multiply the step counts to count complete outcomes.\par}\end{overlayarea}
\end{frame}

% Frame 11. Teaching note: Collect predictions before revealing outcomes. Numbers identify objects; never multiply the identifiers.
\begin{frame}[t]{2A. Arrange both distinct objects}
\Q{Predict: how many different outcomes are possible?}
\DiagramStart
\Obj{C}{4.8}{2.8}{1}\Obj{C}{5.9}{2.8}{2}\node[] at (6.7,1.7) {Available distinct objects};\draw (4.9,.6) rectangle (5.5,1.25);\node[font=\small] at (5.2,0.25) {1};\draw (5.8500000000000005,.6) rectangle (6.45,1.25);\node[font=\small] at (6.15,0.25) {2};\node[font=\small] at (6.75,-0.3) {Empty left-to-right positions};\node[font=\small] at (6.75,1.55) {Use every object once. Read positions from left to right.};
\DiagramEnd
\begin{overlayarea}{\linewidth}{1.15cm}\end{overlayarea}
\end{frame}

\section{2. Arrange all distinct objects}
% Frame 12. Teaching note: Position numbers below the rows are separate from identifiers above the circles.
\begin{frame}[t]{2A. Both possible rows}
\Q{Does exchanging the objects change the row?}
\DiagramStart
\visible<2->{\Obj{C}{5.6}{2.6}{1}\node[font=\small] at (5.6,2.2800000000000002) {1};\Obj{C}{6.14}{2.6}{2}\node[font=\small] at (6.14,2.2800000000000002) {2};}\visible<3->{\Obj{C}{5.6}{1}{2}\node[font=\small] at (5.6,0.6799999999999999) {1};\Obj{C}{6.14}{1}{1}\node[font=\small] at (6.14,0.6799999999999999) {2};}\visible<4->{\draw[group] (5.2,2.08) rectangle (6.6,3.18);}\visible<4->{\draw[group] (5.2,0.48) rectangle (6.6,1.58);}
\DiagramEnd
\begin{overlayarea}{\linewidth}{1.15cm}\only<4>{Two first-position groups; one row in each.\par}
\only<5->{$2\times1=2$ rows. Yes: the objects exchange positions.\par}\end{overlayarea}
\end{frame}

% Frame 13. Teaching note: Collect predictions before revealing outcomes. Numbers identify objects; never multiply the identifiers.
\begin{frame}[t]{2B. Arrange three distinct objects}
\Q{Predict: how many different outcomes are possible?}
\DiagramStart
\Obj{C}{4.8}{2.8}{1}\Obj{C}{5.9}{2.8}{2}\Obj{C}{7.0}{2.8}{3}\node[] at (6.7,1.7) {Available distinct objects};\draw (4.9,.6) rectangle (5.5,1.25);\node[font=\small] at (5.2,0.25) {1};\draw (5.8500000000000005,.6) rectangle (6.45,1.25);\node[font=\small] at (6.15,0.25) {2};\draw (6.8,.6) rectangle (7.3999999999999995,1.25);\node[font=\small] at (7.1,0.25) {3};\node[font=\small] at (6.75,-0.3) {Empty left-to-right positions};\node[font=\small] at (6.75,1.55) {Use every object once in three left-to-right positions.};
\DiagramEnd
\begin{overlayarea}{\linewidth}{1.15cm}\end{overlayarea}
\end{frame}

% Frame 14. Teaching note: List complete outcomes before naming factorial. Do not treat these three columns as three outcomes.
\begin{frame}[t]{2B. Three first-position groups}
\Q{How many rows start with each first object?}
\DiagramStart
\node[font=\small] at (1.94,3.5) {First object: 1};\visible<2->{\Obj{C}{1.4}{2.65}{1}\node[font=\small] at (1.4,2.33) {1};\Obj{C}{1.94}{2.65}{2}\node[font=\small] at (1.94,2.33) {2};\Obj{C}{2.48}{2.65}{3}\node[font=\small] at (2.48,2.33) {3};}\visible<2->{\Obj{C}{1.4}{1.2999999999999998}{1}\node[font=\small] at (1.4,0.9799999999999998) {1};\Obj{C}{1.94}{1.2999999999999998}{3}\node[font=\small] at (1.94,0.9799999999999998) {2};\Obj{C}{2.48}{1.2999999999999998}{2}\node[font=\small] at (2.48,0.9799999999999998) {3};}\visible<5->{\draw[group] (0.9999999999999999,0.75) rectangle (2.9,3.2);}\node[font=\small] at (6.14,3.5) {First object: 2};\visible<3->{\Obj{C}{5.6}{2.65}{2}\node[font=\small] at (5.6,2.33) {1};\Obj{C}{6.14}{2.65}{1}\node[font=\small] at (6.14,2.33) {2};\Obj{C}{6.68}{2.65}{3}\node[font=\small] at (6.68,2.33) {3};}\visible<3->{\Obj{C}{5.6}{1.2999999999999998}{2}\node[font=\small] at (5.6,0.9799999999999998) {1};\Obj{C}{6.14}{1.2999999999999998}{3}\node[font=\small] at (6.14,0.9799999999999998) {2};\Obj{C}{6.68}{1.2999999999999998}{1}\node[font=\small] at (6.68,0.9799999999999998) {3};}\visible<5->{\draw[group] (5.199999999999999,0.75) rectangle (7.1,3.2);}\node[font=\small] at (10.34,3.5) {First object: 3};\visible<4->{\Obj{C}{9.8}{2.65}{3}\node[font=\small] at (9.8,2.33) {1};\Obj{C}{10.34}{2.65}{1}\node[font=\small] at (10.34,2.33) {2};\Obj{C}{10.88}{2.65}{2}\node[font=\small] at (10.88,2.33) {3};}\visible<4->{\Obj{C}{9.8}{1.2999999999999998}{3}\node[font=\small] at (9.8,0.9799999999999998) {1};\Obj{C}{10.34}{1.2999999999999998}{2}\node[font=\small] at (10.34,0.9799999999999998) {2};\Obj{C}{10.88}{1.2999999999999998}{1}\node[font=\small] at (10.88,0.9799999999999998) {3};}\visible<5->{\draw[group] (9.4,0.75) rectangle (11.3,3.2);}
\DiagramEnd
\begin{overlayarea}{\linewidth}{1.15cm}\only<5>{Three first choices. Each group has two orders of the remaining objects.\par}
\only<6->{Three positions give $3\times2\times1=6=3!$ rows.\par}\end{overlayarea}
\end{frame}

% Frame 15. Teaching note: Synchronize used-object crosses with filled boxes. Crosses mean already placed, not removed from the outcome.
\begin{frame}[t]{2B. Follow one path through the positions}
\Q{Why do the choice counts decrease?}
\DiagramStart
\node[anchor=west] at (0.4,3.25) {Available:};\Obj{C}{4.5}{3.1}{1}\Obj{C}{5.9}{3.1}{2}\Obj{C}{7.3}{3.1}{3}\draw (2.55,.9) rectangle (3.45,1.8);\node[font=\small] at (3.0,0.5) {Position 1};\visible<2->{\Obj{C}{3.0}{1.25}{1}}\visible<2->{\draw (4.26,2.88)--(4.74,3.55);\draw (4.74,2.88)--(4.26,3.55);}\visible<2->{\node[font=\small] at (3.0,0) {3 choices};}\draw (5.85,.9) rectangle (6.75,1.8);\node[font=\small] at (6.3,0.5) {Position 2};\visible<3->{\Obj{C}{6.3}{1.25}{2}}\visible<3->{\draw (5.66,2.88)--(6.140000000000001,3.55);\draw (6.140000000000001,2.88)--(5.66,3.55);}\visible<3->{\node[font=\small] at (6.3,0) {2 choices};}\draw (9.15,.9) rectangle (10.049999999999999,1.8);\node[font=\small] at (9.6,0.5) {Position 3};\visible<4->{\Obj{C}{9.6}{1.25}{3}}\visible<4->{\draw (7.06,2.88)--(7.54,3.55);\draw (7.54,2.88)--(7.06,3.55);}\visible<4->{\node[font=\small] at (9.6,0) {1 choice};}
\DiagramEnd
\begin{overlayarea}{\linewidth}{1.15cm}\only<2>{Use object 1: it cannot be chosen again.\par}
\only<3>{Use object 2: only object 3 remains.\par}
\only<4->{Every other path has the same counts: three, then two, then one.\par}\end{overlayarea}
\end{frame}

% Frame 16. Teaching note: Collect predictions before revealing outcomes. Numbers identify objects; never multiply the identifiers.
\begin{frame}[t]{2C. Arrange four distinct objects}
\Q{Predict: how many different outcomes are possible?}
\DiagramStart
\Obj{C}{4.8}{2.8}{1}\Obj{C}{5.9}{2.8}{2}\Obj{C}{7.0}{2.8}{3}\Obj{C}{8.1}{2.8}{4}\node[] at (6.7,1.7) {Available distinct objects};\draw (4.9,.6) rectangle (5.5,1.25);\node[font=\small] at (5.2,0.25) {1};\draw (5.8500000000000005,.6) rectangle (6.45,1.25);\node[font=\small] at (6.15,0.25) {2};\draw (6.8,.6) rectangle (7.3999999999999995,1.25);\node[font=\small] at (7.1,0.25) {3};\draw (7.750000000000001,.6) rectangle (8.350000000000001,1.25);\node[font=\small] at (8.05,0.25) {4};\node[font=\small] at (6.75,-0.3) {Empty left-to-right positions};\node[font=\small] at (6.75,1.55) {Use every object once in four left-to-right positions.};
\DiagramEnd
\begin{overlayarea}{\linewidth}{1.15cm}\end{overlayarea}
\end{frame}

% Frame 17. Teaching note: All twelve rows on this frame are complete; the next frame keeps identifiers and shows the other two blocks.
\begin{frame}[t]{2C. Complete rows starting with 1 or 2}
\Q{How many second-position subgroups fit inside each block?}
\DiagramStart
\node[font=\small] at (2.0,3.85) {First object: 1};\visible<2->{\Obj{C}{1.0}{3.2}{1}\Obj{C}{1.6}{3.2}{2}\Obj{C}{2.2}{3.2}{3}\Obj{C}{2.8}{3.2}{4}}\visible<2->{\Obj{C}{1.0}{2.5500000000000003}{1}\Obj{C}{1.6}{2.5500000000000003}{2}\Obj{C}{2.2}{2.5500000000000003}{4}\Obj{C}{2.8}{2.5500000000000003}{3}}\visible<5->{\draw[group] (0.65,2.24) rectangle (3.2,3.68);}\visible<5->{\node[font=\small] at (4.55,3.1) {Second: 2};}\visible<3->{\Obj{C}{1.0}{1.9000000000000001}{1}\Obj{C}{1.6}{1.9000000000000001}{3}\Obj{C}{2.2}{1.9000000000000001}{2}\Obj{C}{2.8}{1.9000000000000001}{4}}\visible<3->{\Obj{C}{1.0}{1.25}{1}\Obj{C}{1.6}{1.25}{3}\Obj{C}{2.2}{1.25}{4}\Obj{C}{2.8}{1.25}{2}}\visible<5->{\draw[group] (0.65,0.9400000000000002) rectangle (3.2,2.38);}\visible<5->{\node[font=\small] at (4.55,1.8) {Second: 3};}\visible<4->{\Obj{C}{1.0}{0.6000000000000001}{1}\Obj{C}{1.6}{0.6000000000000001}{4}\Obj{C}{2.2}{0.6000000000000001}{2}\Obj{C}{2.8}{0.6000000000000001}{3}}\visible<4->{\Obj{C}{1.0}{-0.04999999999999993}{1}\Obj{C}{1.6}{-0.04999999999999993}{4}\Obj{C}{2.2}{-0.04999999999999993}{3}\Obj{C}{2.8}{-0.04999999999999993}{2}}\visible<5->{\draw[group] (0.65,-0.3599999999999999) rectangle (3.2,1.08);}\visible<5->{\node[font=\small] at (4.55,0.5000000000000001) {Second: 4};}\visible<6->{\draw[group] (0.44999999999999996,-0.5) rectangle (5.65,4.05);}\node[font=\small] at (8.5,3.85) {First object: 2};\visible<2->{\Obj{C}{7.5}{3.2}{2}\Obj{C}{8.1}{3.2}{1}\Obj{C}{8.7}{3.2}{3}\Obj{C}{9.3}{3.2}{4}}\visible<2->{\Obj{C}{7.5}{2.5500000000000003}{2}\Obj{C}{8.1}{2.5500000000000003}{1}\Obj{C}{8.7}{2.5500000000000003}{4}\Obj{C}{9.3}{2.5500000000000003}{3}}\visible<5->{\draw[group] (7.15,2.24) rectangle (9.7,3.68);}\visible<5->{\node[font=\small] at (11.05,3.1) {Second: 1};}\visible<3->{\Obj{C}{7.5}{1.9000000000000001}{2}\Obj{C}{8.1}{1.9000000000000001}{3}\Obj{C}{8.7}{1.9000000000000001}{1}\Obj{C}{9.3}{1.9000000000000001}{4}}\visible<3->{\Obj{C}{7.5}{1.25}{2}\Obj{C}{8.1}{1.25}{3}\Obj{C}{8.7}{1.25}{4}\Obj{C}{9.3}{1.25}{1}}\visible<5->{\draw[group] (7.15,0.9400000000000002) rectangle (9.7,2.38);}\visible<5->{\node[font=\small] at (11.05,1.8) {Second: 3};}\visible<4->{\Obj{C}{7.5}{0.6000000000000001}{2}\Obj{C}{8.1}{0.6000000000000001}{4}\Obj{C}{8.7}{0.6000000000000001}{1}\Obj{C}{9.3}{0.6000000000000001}{3}}\visible<4->{\Obj{C}{7.5}{-0.04999999999999993}{2}\Obj{C}{8.1}{-0.04999999999999993}{4}\Obj{C}{8.7}{-0.04999999999999993}{3}\Obj{C}{9.3}{-0.04999999999999993}{1}}\visible<5->{\draw[group] (7.15,-0.3599999999999999) rectangle (9.7,1.08);}\visible<5->{\node[font=\small] at (11.05,0.5000000000000001) {Second: 4};}\visible<6->{\draw[group] (6.95,-0.5) rectangle (12.15,4.05);}
\DiagramEnd
\begin{overlayarea}{\linewidth}{1.15cm}\only<5>{Three second-position subgroups; two rows in every subgroup.\par}
\only<6->{Each first-position block has six complete rows. The last object is forced.\par}\end{overlayarea}
\end{frame}

% Frame 18. Teaching note: All twelve rows on this frame are complete; the next frame keeps identifiers and shows the other two blocks.
\begin{frame}[t]{2C. Complete rows starting with 3 or 4}
\Q{How many second-position subgroups fit inside each block?}
\DiagramStart
\node[font=\small] at (2.0,3.85) {First object: 3};\visible<2->{\Obj{C}{1.0}{3.2}{3}\Obj{C}{1.6}{3.2}{1}\Obj{C}{2.2}{3.2}{2}\Obj{C}{2.8}{3.2}{4}}\visible<2->{\Obj{C}{1.0}{2.5500000000000003}{3}\Obj{C}{1.6}{2.5500000000000003}{1}\Obj{C}{2.2}{2.5500000000000003}{4}\Obj{C}{2.8}{2.5500000000000003}{2}}\visible<5->{\draw[group] (0.65,2.24) rectangle (3.2,3.68);}\visible<5->{\node[font=\small] at (4.55,3.1) {Second: 1};}\visible<3->{\Obj{C}{1.0}{1.9000000000000001}{3}\Obj{C}{1.6}{1.9000000000000001}{2}\Obj{C}{2.2}{1.9000000000000001}{1}\Obj{C}{2.8}{1.9000000000000001}{4}}\visible<3->{\Obj{C}{1.0}{1.25}{3}\Obj{C}{1.6}{1.25}{2}\Obj{C}{2.2}{1.25}{4}\Obj{C}{2.8}{1.25}{1}}\visible<5->{\draw[group] (0.65,0.9400000000000002) rectangle (3.2,2.38);}\visible<5->{\node[font=\small] at (4.55,1.8) {Second: 2};}\visible<4->{\Obj{C}{1.0}{0.6000000000000001}{3}\Obj{C}{1.6}{0.6000000000000001}{4}\Obj{C}{2.2}{0.6000000000000001}{1}\Obj{C}{2.8}{0.6000000000000001}{2}}\visible<4->{\Obj{C}{1.0}{-0.04999999999999993}{3}\Obj{C}{1.6}{-0.04999999999999993}{4}\Obj{C}{2.2}{-0.04999999999999993}{2}\Obj{C}{2.8}{-0.04999999999999993}{1}}\visible<5->{\draw[group] (0.65,-0.3599999999999999) rectangle (3.2,1.08);}\visible<5->{\node[font=\small] at (4.55,0.5000000000000001) {Second: 4};}\visible<6->{\draw[group] (0.44999999999999996,-0.5) rectangle (5.65,4.05);}\node[font=\small] at (8.5,3.85) {First object: 4};\visible<2->{\Obj{C}{7.5}{3.2}{4}\Obj{C}{8.1}{3.2}{1}\Obj{C}{8.7}{3.2}{2}\Obj{C}{9.3}{3.2}{3}}\visible<2->{\Obj{C}{7.5}{2.5500000000000003}{4}\Obj{C}{8.1}{2.5500000000000003}{1}\Obj{C}{8.7}{2.5500000000000003}{3}\Obj{C}{9.3}{2.5500000000000003}{2}}\visible<5->{\draw[group] (7.15,2.24) rectangle (9.7,3.68);}\visible<5->{\node[font=\small] at (11.05,3.1) {Second: 1};}\visible<3->{\Obj{C}{7.5}{1.9000000000000001}{4}\Obj{C}{8.1}{1.9000000000000001}{2}\Obj{C}{8.7}{1.9000000000000001}{1}\Obj{C}{9.3}{1.9000000000000001}{3}}\visible<3->{\Obj{C}{7.5}{1.25}{4}\Obj{C}{8.1}{1.25}{2}\Obj{C}{8.7}{1.25}{3}\Obj{C}{9.3}{1.25}{1}}\visible<5->{\draw[group] (7.15,0.9400000000000002) rectangle (9.7,2.38);}\visible<5->{\node[font=\small] at (11.05,1.8) {Second: 2};}\visible<4->{\Obj{C}{7.5}{0.6000000000000001}{4}\Obj{C}{8.1}{0.6000000000000001}{3}\Obj{C}{8.7}{0.6000000000000001}{1}\Obj{C}{9.3}{0.6000000000000001}{2}}\visible<4->{\Obj{C}{7.5}{-0.04999999999999993}{4}\Obj{C}{8.1}{-0.04999999999999993}{3}\Obj{C}{8.7}{-0.04999999999999993}{2}\Obj{C}{9.3}{-0.04999999999999993}{1}}\visible<5->{\draw[group] (7.15,-0.3599999999999999) rectangle (9.7,1.08);}\visible<5->{\node[font=\small] at (11.05,0.5000000000000001) {Second: 3};}\visible<6->{\draw[group] (6.95,-0.5) rectangle (12.15,4.05);}
\DiagramEnd
\begin{overlayarea}{\linewidth}{1.15cm}\only<5>{Three second-position subgroups; two rows in every subgroup.\par}
\only<6->{Each first-position block has six complete rows. The last object is forced.\par}\end{overlayarea}
\end{frame}

% Frame 19. Teaching note: The four six-row blocks correspond exactly to the preceding two frames; none is omitted.
\begin{frame}[t]{2C. Put the four blocks together}
\Q{What does each factor count?}
\DiagramStart
\visible<2->{\draw[group] (0.6000000000000001,1) rectangle (2.6,3.4);}\Obj{C}{1.6}{2.8}{1}\visible<2->{\node[] at (1.6,1.7) {6 rows};}\visible<2->{\draw[group] (3.95,1) rectangle (5.95,3.4);}\Obj{C}{4.95}{2.8}{2}\visible<2->{\node[] at (4.95,1.7) {6 rows};}\visible<2->{\draw[group] (7.300000000000001,1) rectangle (9.3,3.4);}\Obj{C}{8.3}{2.8}{3}\visible<2->{\node[] at (8.3,1.7) {6 rows};}\visible<2->{\draw[group] (10.65,1) rectangle (12.65,3.4);}\Obj{C}{11.65}{2.8}{4}\visible<2->{\node[] at (11.65,1.7) {6 rows};}
\DiagramEnd
\begin{overlayarea}{\linewidth}{1.15cm}\only<2>{Four first-position blocks, with six rows in each.\par}
\only<3>{$4\times3\times2\times1=24=4!$.\par}
\only<4->{First choice; second choice; third choice; one last object.\par}\end{overlayarea}
\end{frame}

% Frame 20. Teaching note: Factorial is introduced after all three concrete examples. The exclamation mark is an operation, not another factor.
\begin{frame}[t]{Factorial names a decreasing product}
\Q{How would this work for any number of distinct objects?}
\DiagramStart
\visible<3->{\node[font=\Large] at (6.75,2.7) {$n!=n\times(n-1)\times\cdots\times2\times1$};}\visible<4->{\node[] at (6.75,1.2) {$4!=4\times3\times2\times1$};}
\DiagramEnd
\begin{overlayarea}{\linewidth}{1.15cm}\only<2>{$n$ is the total number of objects, a whole number.\par}
\only<3>{$!$ means factorial: multiply down to 1. Here $n\geq1$.\par}
\only<4->{$n-1$ means one fewer; the dots mean keep decreasing by one.\par}\end{overlayarea}
\end{frame}

% Frame 21. Teaching note: Zero objects is not zero outcomes: the one empty arrangement is the outcome.
\begin{frame}[t]{One object, and the empty arrangement}
\Q{How many arrangements can one object make? How about none?}
\DiagramStart
\Obj{C}{3.0}{2.5}{1}\visible<2->{\node[] at (3,1.2) {$1!=1$};}\node[] at (10,2.5) {[ empty row ]};\visible<3->{\node[] at (10,1.2) {$0!=1$};}
\DiagramEnd
\begin{overlayarea}{\linewidth}{1.15cm}\only<2>{One object has exactly one arrangement.\par}
\only<3->{No objects have one empty arrangement. We define zero factorial as 1.\par}\end{overlayarea}
\end{frame}

% Frame 22. Teaching note: Collect predictions before revealing outcomes. Numbers identify objects; never multiply the identifiers.
\begin{frame}[t]{3A. One Reader from three objects}
\Q{Predict: how many different outcomes are possible?}
\DiagramStart
\Obj{C}{4.8}{2.8}{1}\Obj{C}{5.9}{2.8}{2}\Obj{C}{7.0}{2.8}{3}\node[] at (6.7,1.7) {Available distinct objects};\node[] at (6.75,0.5) {One assignment contains only the selected Reader.};
\DiagramEnd
\begin{overlayarea}{\linewidth}{1.15cm}\end{overlayarea}
\end{frame}

\section{3. Fill some ordered positions}
% Frame 23. Teaching note: Unused objects are absent from the assignment. Only the selected Reader is counted.
\begin{frame}[t]{3A. Three one-position assignments}
\Q{How many choices fill the single job?}
\DiagramStart
\visible<2->{\Obj{C}{2.0}{2.6}{1}}\node[font=\small] at (2.0,1.9) {Reader};\visible<3->{\Obj{C}{6.5}{2.6}{2}}\node[font=\small] at (6.5,1.9) {Reader};\visible<4->{\Obj{C}{11.0}{2.6}{3}}\node[font=\small] at (11.0,1.9) {Reader};
\DiagramEnd
\begin{overlayarea}{\linewidth}{1.15cm}\only<5->{Three choices for one position: three assignments.\par}\end{overlayarea}
\end{frame}

% Frame 24. Teaching note: The divisor counts orders of unused objects, never the selected Reader. Keep faded tails visible.
\begin{frame}[t]{3A. Full rows overcount the Reader assignment}
\Q{Do the two faded tail orders change the Reader?}
\DiagramStart
\visible<2->{\Obj{C}{1.5}{2.7}{1}\begin{scope}[opacity=.30]\Obj{C}{2.15}{2.7}{2}\end{scope}\begin{scope}[opacity=.30]\Obj{C}{2.8}{2.7}{3}\end{scope}}\visible<2->{\Obj{C}{1.5}{1.4000000000000001}{1}\begin{scope}[opacity=.30]\Obj{C}{2.15}{1.4000000000000001}{3}\end{scope}\begin{scope}[opacity=.30]\Obj{C}{2.8}{1.4000000000000001}{2}\end{scope}}\node[font=\small] at (1.5,0.55) {Reader};\node[font=\small] at (2.5,0.55) {unused tail};\visible<3->{\draw[group] (1.1,0.9) rectangle (3.2,3.3);}\visible<2->{\Obj{C}{5.7}{2.7}{2}\begin{scope}[opacity=.30]\Obj{C}{6.35}{2.7}{1}\end{scope}\begin{scope}[opacity=.30]\Obj{C}{7.0}{2.7}{3}\end{scope}}\visible<2->{\Obj{C}{5.7}{1.4000000000000001}{2}\begin{scope}[opacity=.30]\Obj{C}{6.35}{1.4000000000000001}{3}\end{scope}\begin{scope}[opacity=.30]\Obj{C}{7.0}{1.4000000000000001}{1}\end{scope}}\node[font=\small] at (5.7,0.55) {Reader};\node[font=\small] at (6.7,0.55) {unused tail};\visible<3->{\draw[group] (5.3,0.9) rectangle (7.4,3.3);}\visible<2->{\Obj{C}{9.9}{2.7}{3}\begin{scope}[opacity=.30]\Obj{C}{10.55}{2.7}{1}\end{scope}\begin{scope}[opacity=.30]\Obj{C}{11.2}{2.7}{2}\end{scope}}\visible<2->{\Obj{C}{9.9}{1.4000000000000001}{3}\begin{scope}[opacity=.30]\Obj{C}{10.55}{1.4000000000000001}{2}\end{scope}\begin{scope}[opacity=.30]\Obj{C}{11.2}{1.4000000000000001}{1}\end{scope}}\node[font=\small] at (9.9,0.55) {Reader};\node[font=\small] at (10.9,0.55) {unused tail};\visible<3->{\draw[group] (9.5,0.9) rectangle (11.6,3.3);}
\DiagramEnd
\begin{overlayarea}{\linewidth}{1.15cm}\only<3>{No. Both rows in a box give the same one-position assignment.\par}
\only<4>{Each assignment has $2!$ orders of its unused objects.\par}
\only<5->{$3!\div2!=6\div2=3$ Reader assignments.\par}\end{overlayarea}
\end{frame}

% Frame 25. Teaching note: Collect predictions before revealing outcomes. Numbers identify objects; never multiply the identifiers.
\begin{frame}[t]{3B. Reader and Recorder from three}
\Q{Predict: how many different outcomes are possible?}
\DiagramStart
\Obj{C}{4.8}{2.8}{1}\Obj{C}{5.9}{2.8}{2}\Obj{C}{7.0}{2.8}{3}\node[] at (6.7,1.7) {Available distinct objects};\node[] at (6.75,0.5) {One assignment: Reader first, Recorder second. No object repeats.};
\DiagramEnd
\begin{overlayarea}{\linewidth}{1.15cm}\end{overlayarea}
\end{frame}

% Frame 26. Teaching note: Columns fix the Reader. Recorder excludes that object. No repeated-object rows.
\begin{frame}[t]{3B. Every ordered two-job assignment}
\Q{How many choices remain after the first object is used?}
\DiagramStart
\node[font=\small] at (1.15,3.55) {First: 1};\visible<2->{\Obj{C}{0.65}{2.75}{1}\Obj{C}{1.9}{2.75}{2}}\visible<2->{\node[font=\small] at (0.65,2.43) {Reader};}\visible<2->{\node[font=\small] at (1.9,2.43) {Recorder};}\visible<2->{\Obj{C}{0.65}{1.85}{1}\Obj{C}{1.9}{1.85}{3}}\visible<2->{\node[font=\small] at (0.65,1.53) {Reader};}\visible<2->{\node[font=\small] at (1.9,1.53) {Recorder};}\visible<5->{\draw[group] (0.0,1.33) rectangle (2.7,3.3);}\node[font=\small] at (5.416666666666667,3.55) {First: 2};\visible<3->{\Obj{C}{4.917}{2.75}{2}\Obj{C}{6.167}{2.75}{1}}\visible<3->{\node[font=\small] at (4.916666666666667,2.43) {Reader};}\visible<3->{\node[font=\small] at (6.166666666666667,2.43) {Recorder};}\visible<3->{\Obj{C}{4.917}{1.85}{2}\Obj{C}{6.167}{1.85}{3}}\visible<3->{\node[font=\small] at (4.916666666666667,1.53) {Reader};}\visible<3->{\node[font=\small] at (6.166666666666667,1.53) {Recorder};}\visible<5->{\draw[group] (4.266666666666667,1.33) rectangle (6.966666666666667,3.3);}\node[font=\small] at (9.683333333333334,3.55) {First: 3};\visible<4->{\Obj{C}{9.183}{2.75}{3}\Obj{C}{10.433}{2.75}{1}}\visible<4->{\node[font=\small] at (9.183333333333334,2.43) {Reader};}\visible<4->{\node[font=\small] at (10.433333333333334,2.43) {Recorder};}\visible<4->{\Obj{C}{9.183}{1.85}{3}\Obj{C}{10.433}{1.85}{2}}\visible<4->{\node[font=\small] at (9.183333333333334,1.53) {Reader};}\visible<4->{\node[font=\small] at (10.433333333333334,1.53) {Recorder};}\visible<5->{\draw[group] (8.533333333333333,1.33) rectangle (11.233333333333334,3.3);}
\DiagramEnd
\begin{overlayarea}{\linewidth}{1.15cm}\only<5>{3 first-object groups; 2 assignments in each.\par}
\only<6->{$3\times2=6$ assignments. Each object is used at most once.\par}\end{overlayarea}
\end{frame}

% Frame 27. Teaching note: Every two-job assignment has exactly one full-row completion; do not divide by two job orders.
\begin{frame}[t]{3B. The unused object has one order}
\Q{Why is the unused-tail divisor now only one?}
\DiagramStart
\Obj{C}{5.0}{2.7}{1}\Obj{C}{6.2}{2.7}{2}\begin{scope}[opacity=.30]\Obj{C}{7.4}{2.7}{3}\end{scope}\node[font=\small] at (5,1.8) {Reader};\node[font=\small] at (6.2,1.8) {Recorder};\node[font=\small] at (7.4,1.8) {Unused};
\DiagramEnd
\begin{overlayarea}{\linewidth}{1.15cm}\only<2>{After two jobs are filled, only one object is unused.\par}
\only<3->{Its single order contributes $1!$: $3!\div1!=6$ assignments.\par}\end{overlayarea}
\end{frame}

% Frame 28. Teaching note: Collect predictions before revealing outcomes. Numbers identify objects; never multiply the identifiers.
\begin{frame}[t]{3C. Reader and Recorder from four}
\Q{Predict: how many different outcomes are possible?}
\DiagramStart
\Obj{C}{4.8}{2.8}{1}\Obj{C}{5.9}{2.8}{2}\Obj{C}{7.0}{2.8}{3}\Obj{C}{8.1}{2.8}{4}\node[] at (6.7,1.7) {Available distinct objects};\node[] at (6.75,0.5) {One assignment has two job holders. The other two are unused.};
\DiagramEnd
\begin{overlayarea}{\linewidth}{1.15cm}\end{overlayarea}
\end{frame}

% Frame 29. Teaching note: Columns fix the Reader. Recorder excludes that object. No repeated-object rows.
\begin{frame}[t]{3C. Four columns of three assignments}
\Q{How many choices remain after the first object is used?}
\DiagramStart
\node[font=\small] at (1.15,3.55) {First: 1};\visible<2->{\Obj{C}{0.65}{2.75}{1}\Obj{C}{1.9}{2.75}{2}}\visible<2->{\node[font=\small] at (0.65,2.43) {Reader};}\visible<2->{\node[font=\small] at (1.9,2.43) {Recorder};}\visible<2->{\Obj{C}{0.65}{1.85}{1}\Obj{C}{1.9}{1.85}{3}}\visible<2->{\node[font=\small] at (0.65,1.53) {Reader};}\visible<2->{\node[font=\small] at (1.9,1.53) {Recorder};}\visible<2->{\Obj{C}{0.65}{0.95}{1}\Obj{C}{1.9}{0.95}{4}}\visible<2->{\node[font=\small] at (0.65,0.6299999999999999) {Reader};}\visible<2->{\node[font=\small] at (1.9,0.6299999999999999) {Recorder};}\visible<6->{\draw[group] (0.0,0.42999999999999994) rectangle (2.7,3.3);}\node[font=\small] at (4.35,3.55) {First: 2};\visible<3->{\Obj{C}{3.85}{2.75}{2}\Obj{C}{5.1}{2.75}{1}}\visible<3->{\node[font=\small] at (3.85,2.43) {Reader};}\visible<3->{\node[font=\small] at (5.1,2.43) {Recorder};}\visible<3->{\Obj{C}{3.85}{1.85}{2}\Obj{C}{5.1}{1.85}{3}}\visible<3->{\node[font=\small] at (3.85,1.53) {Reader};}\visible<3->{\node[font=\small] at (5.1,1.53) {Recorder};}\visible<3->{\Obj{C}{3.85}{0.95}{2}\Obj{C}{5.1}{0.95}{4}}\visible<3->{\node[font=\small] at (3.85,0.6299999999999999) {Reader};}\visible<3->{\node[font=\small] at (5.1,0.6299999999999999) {Recorder};}\visible<6->{\draw[group] (3.2,0.42999999999999994) rectangle (5.8999999999999995,3.3);}\node[font=\small] at (7.550000000000001,3.55) {First: 3};\visible<4->{\Obj{C}{7.05}{2.75}{3}\Obj{C}{8.3}{2.75}{1}}\visible<4->{\node[font=\small] at (7.050000000000001,2.43) {Reader};}\visible<4->{\node[font=\small] at (8.3,2.43) {Recorder};}\visible<4->{\Obj{C}{7.05}{1.85}{3}\Obj{C}{8.3}{1.85}{2}}\visible<4->{\node[font=\small] at (7.050000000000001,1.53) {Reader};}\visible<4->{\node[font=\small] at (8.3,1.53) {Recorder};}\visible<4->{\Obj{C}{7.05}{0.95}{3}\Obj{C}{8.3}{0.95}{4}}\visible<4->{\node[font=\small] at (7.050000000000001,0.6299999999999999) {Reader};}\visible<4->{\node[font=\small] at (8.3,0.6299999999999999) {Recorder};}\visible<6->{\draw[group] (6.4,0.42999999999999994) rectangle (9.100000000000001,3.3);}\node[font=\small] at (10.750000000000002,3.55) {First: 4};\visible<5->{\Obj{C}{10.25}{2.75}{4}\Obj{C}{11.5}{2.75}{1}}\visible<5->{\node[font=\small] at (10.250000000000002,2.43) {Reader};}\visible<5->{\node[font=\small] at (11.500000000000002,2.43) {Recorder};}\visible<5->{\Obj{C}{10.25}{1.85}{4}\Obj{C}{11.5}{1.85}{2}}\visible<5->{\node[font=\small] at (10.250000000000002,1.53) {Reader};}\visible<5->{\node[font=\small] at (11.500000000000002,1.53) {Recorder};}\visible<5->{\Obj{C}{10.25}{0.95}{4}\Obj{C}{11.5}{0.95}{3}}\visible<5->{\node[font=\small] at (10.250000000000002,0.6299999999999999) {Reader};}\visible<5->{\node[font=\small] at (11.500000000000002,0.6299999999999999) {Recorder};}\visible<6->{\draw[group] (9.600000000000001,0.42999999999999994) rectangle (12.300000000000002,3.3);}
\DiagramEnd
\begin{overlayarea}{\linewidth}{1.15cm}\only<6>{4 first-object groups; 3 assignments in each.\par}
\only<7->{$4\times3=12$ assignments. Each object is used at most once.\par}\end{overlayarea}
\end{frame}

% Frame 30. Teaching note: Left grid remains complete. Right boxes illustrate 12 and 13; every one of the twelve has exactly the same two-completion structure.
\begin{frame}[t]{3C. Completion boxes explain the divisor}
\Q{How many full rows complete each two-job assignment?}
\DiagramStart
\Obj{C}{0.5}{2.9}{1}\Obj{C}{1.0}{2.9}{2}\Obj{C}{0.5}{1.95}{1}\Obj{C}{1.0}{1.95}{3}\Obj{C}{0.5}{1.0}{1}\Obj{C}{1.0}{1.0}{4}\Obj{C}{2.15}{2.9}{2}\Obj{C}{2.65}{2.9}{1}\Obj{C}{2.15}{1.95}{2}\Obj{C}{2.65}{1.95}{3}\Obj{C}{2.15}{1.0}{2}\Obj{C}{2.65}{1.0}{4}\Obj{C}{3.8}{2.9}{3}\Obj{C}{4.3}{2.9}{1}\Obj{C}{3.8}{1.95}{3}\Obj{C}{4.3}{1.95}{2}\Obj{C}{3.8}{1.0}{3}\Obj{C}{4.3}{1.0}{4}\Obj{C}{5.45}{2.9}{4}\Obj{C}{5.95}{2.9}{1}\Obj{C}{5.45}{1.95}{4}\Obj{C}{5.95}{1.95}{2}\Obj{C}{5.45}{1.0}{4}\Obj{C}{5.95}{1.0}{3}\node[font=\small] at (3.1,0.03) {All twelve assignments};\visible<2->{\Obj{C}{8.1}{3.0}{1}\Obj{C}{8.9}{3.0}{2}\begin{scope}[opacity=.30]\Obj{C}{9.7}{3.0}{3}\end{scope}\begin{scope}[opacity=.30]\Obj{C}{10.5}{3.0}{4}\end{scope}}\visible<2->{\Obj{C}{8.1}{2.35}{1}\Obj{C}{8.9}{2.35}{2}\begin{scope}[opacity=.30]\Obj{C}{9.7}{2.35}{4}\end{scope}\begin{scope}[opacity=.30]\Obj{C}{10.5}{2.35}{3}\end{scope}}\visible<2->{\draw[group] (7.8,2.06) rectangle (10.899999999999999,3.55);}\visible<2->{\draw[decorate,decoration={brace,mirror}] (9.45,2.0)--(10.8,2.0);}\visible<2->{\node[font=\small] at (10.1,1.75) {unused tail};}\visible<3->{\Obj{C}{8.1}{1.2}{1}\Obj{C}{8.9}{1.2}{3}\begin{scope}[opacity=.30]\Obj{C}{9.7}{1.2}{2}\end{scope}\begin{scope}[opacity=.30]\Obj{C}{10.5}{1.2}{4}\end{scope}}\visible<3->{\Obj{C}{8.1}{0.5499999999999999}{1}\Obj{C}{8.9}{0.5499999999999999}{3}\begin{scope}[opacity=.30]\Obj{C}{9.7}{0.5499999999999999}{4}\end{scope}\begin{scope}[opacity=.30]\Obj{C}{10.5}{0.5499999999999999}{2}\end{scope}}\visible<3->{\draw[group] (7.8,0.26) rectangle (10.899999999999999,1.75);}\visible<3->{\draw[decorate,decoration={brace,mirror}] (9.45,0.19999999999999996)--(10.8,0.19999999999999996);}\visible<3->{\node[font=\small] at (10.1,-0.050000000000000044) {unused tail};}
\DiagramEnd
\begin{overlayarea}{\linewidth}{1.15cm}\only<4>{Every assignment leaves two unused objects: exactly two tail orders.\par}
\only<5->{$24$ full rows $\div2$ completions per assignment $=12$ assignments.\par}\end{overlayarea}
\end{frame}

% Frame 31. Teaching note: Cancellation removes equal multiplicities. It does not identify Reader--Recorder with Recorder--Reader.
\begin{frame}[t]{3C. Cancel the unused tail}
\Q{Which factors count only unused-tail orders?}
\DiagramStart
\visible<2->{\node[] at (6.75,2.8) {$\displaystyle\frac{4!}{2!}=\frac{4\times3\times\underbrace{(2\times1)}_{\text{unused orders}}}{2\times1}$};}\visible<3->{\node[] at (6.75,0.9) {$\displaystyle\frac{4\times3\times\cancel{(2\times1)}}{\cancel{(2\times1)}}=4\times3=12$};}
\DiagramEnd
\begin{overlayarea}{\linewidth}{1.15cm}\only<2>{The selected positions keep their order.\par}
\only<3->{Divide by the two orders of the unused tail.\par}\end{overlayarea}
\end{frame}

% Frame 32. Teaching note: For r=0 there is one empty assignment; for r=n the divisor is 0!=1. The dots represent r factors, even if r is small.
\begin{frame}[t]{Some ordered positions: the general rule}
\Q{What do the symbols count?}
\DiagramStart
\visible<2->{\node[] at (6.75,2.9) {$n$ available; $r$ positions; $n-r$ unused};}\visible<3->{\node[] at (6.75,1.5) {$\displaystyle n(n-1)\cdots(n-r+1)=\frac{n!}{(n-r)!}$};}\visible<3->{\node[] at (6.75,0.3) {Keep $r$ decreasing choice counts.};}
\DiagramEnd
\begin{overlayarea}{\linewidth}{1.15cm}\only<2>{$n,r$ are whole numbers; $0\leq r\leq n$. No object is used twice.\par}
\only<3>{The last position has $n-r+1$ choices: $r-1$ objects are already used.\par}
\only<4->{Or divide full rows by the $(n-r)!$ orders of each unused tail.\par}\end{overlayarea}
\end{frame}

% Frame 33. Teaching note: Collect predictions before revealing outcomes. Numbers identify objects; never multiply the identifiers.
\begin{frame}[t]{4A. Choose two members from three}
\Q{Predict: how many different outcomes are possible?}
\DiagramStart
\Obj{C}{4.8}{2.8}{1}\Obj{C}{5.9}{2.8}{2}\Obj{C}{7.0}{2.8}{3}\node[] at (6.7,1.7) {Available distinct objects};\node[] at (6.75,0.5) {No jobs or position labels. A group contains members, not roles.};
\DiagramEnd
\begin{overlayarea}{\linewidth}{1.15cm}\end{overlayarea}
\end{frame}

\section{4. Choose a group without roles}
% Frame 34. Teaching note: Each box pairs ab with ba. The oval below represents that single membership, not an extra outcome.
\begin{frame}[t]{4A. Six lists become three groups}
\Q{How many ordered lists name the same group?}
\DiagramStart
\visible<2->{\Obj{C}{1.2}{3.0}{1}\Obj{C}{1.74}{3.0}{2}}\visible<2->{\Obj{C}{2.8}{3.0}{2}\Obj{C}{3.34}{3.0}{1}}\visible<3->{\draw[group] (0.85,2.65) rectangle (3.7,3.55);}\visible<4->{\draw (2.2800000000000002,2.15) ellipse (0.73 and .48);\Obj{C}{2.0}{2.05}{1}\Obj{C}{2.56}{2.05}{2}}\visible<2->{\Obj{C}{5.6}{3.0}{1}\Obj{C}{6.14}{3.0}{3}}\visible<2->{\Obj{C}{7.2}{3.0}{3}\Obj{C}{7.74}{3.0}{1}}\visible<3->{\draw[group] (5.250000000000001,2.65) rectangle (8.100000000000001,3.55);}\visible<4->{\draw (6.680000000000001,2.15) ellipse (0.73 and .48);\Obj{C}{6.4}{2.05}{1}\Obj{C}{6.96}{2.05}{3}}\visible<2->{\Obj{C}{10.0}{3.0}{2}\Obj{C}{10.54}{3.0}{3}}\visible<2->{\Obj{C}{11.6}{3.0}{3}\Obj{C}{12.14}{3.0}{2}}\visible<3->{\draw[group] (9.65,2.65) rectangle (12.5,3.55);}\visible<4->{\draw (11.08,2.15) ellipse (0.73 and .48);\Obj{C}{10.8}{2.05}{2}\Obj{C}{11.36}{2.05}{3}}
\DiagramEnd
\begin{overlayarea}{\linewidth}{1.15cm}\only<3>{6 ordered lists; two lists in each membership box.\par}
\only<4>{3 groups. An oval has members, with no labelled positions.\par}
\only<5->{$6\div2!=3$ distinct groups.\par}\end{overlayarea}
\end{frame}

% Frame 35. Teaching note: Collect predictions before revealing outcomes. Numbers identify objects; never multiply the identifiers.
\begin{frame}[t]{4B. Choose two members from four}
\Q{Predict: how many different outcomes are possible?}
\DiagramStart
\Obj{C}{4.8}{2.8}{1}\Obj{C}{5.9}{2.8}{2}\Obj{C}{7.0}{2.8}{3}\Obj{C}{8.1}{2.8}{4}\node[] at (6.7,1.7) {Available distinct objects};\node[] at (6.75,0.5) {Same objects as before. Remove jobs; keep only membership.};
\DiagramEnd
\begin{overlayarea}{\linewidth}{1.15cm}\end{overlayarea}
\end{frame}

% Frame 36. Teaching note: This is the preceding twelve-assignment grid with job labels removed. Next frame changes arrangement explicitly.
\begin{frame}[t]{4B. Begin with the same twelve ordered lists}
\Q{Which lists contain exactly the same members?}
\DiagramStart
\visible<2->{\Obj{C}{1.0}{2.9}{1}\Obj{C}{1.54}{2.9}{2}}\visible<2->{\Obj{C}{1.0}{1.8499999999999999}{1}\Obj{C}{1.54}{1.8499999999999999}{3}}\visible<2->{\Obj{C}{1.0}{0.7999999999999998}{1}\Obj{C}{1.54}{0.7999999999999998}{4}}\visible<3->{\Obj{C}{4.4}{2.9}{2}\Obj{C}{4.94}{2.9}{1}}\visible<3->{\Obj{C}{4.4}{1.8499999999999999}{2}\Obj{C}{4.94}{1.8499999999999999}{3}}\visible<3->{\Obj{C}{4.4}{0.7999999999999998}{2}\Obj{C}{4.94}{0.7999999999999998}{4}}\visible<4->{\Obj{C}{7.8}{2.9}{3}\Obj{C}{8.34}{2.9}{1}}\visible<4->{\Obj{C}{7.8}{1.8499999999999999}{3}\Obj{C}{8.34}{1.8499999999999999}{2}}\visible<4->{\Obj{C}{7.8}{0.7999999999999998}{3}\Obj{C}{8.34}{0.7999999999999998}{4}}\visible<5->{\Obj{C}{11.2}{2.9}{4}\Obj{C}{11.74}{2.9}{1}}\visible<5->{\Obj{C}{11.2}{1.8499999999999999}{4}\Obj{C}{11.74}{1.8499999999999999}{2}}\visible<5->{\Obj{C}{11.2}{0.7999999999999998}{4}\Obj{C}{11.74}{0.7999999999999998}{3}}
\DiagramEnd
\begin{overlayarea}{\linewidth}{1.15cm}\only<6->{Now reorganize these twelve lists into pairs with the same members.\par}\end{overlayarea}
\end{frame}

% Frame 37. Teaching note: Each box pairs ab with ba. The oval below represents that single membership, not an extra outcome.
\begin{frame}[t]{4B. Pair lists by membership}
\Q{How many ordered lists name the same group?}
\DiagramStart
\visible<2->{\Obj{C}{1.2}{3.0}{1}\Obj{C}{1.74}{3.0}{2}}\visible<2->{\Obj{C}{2.8}{3.0}{2}\Obj{C}{3.34}{3.0}{1}}\visible<3->{\draw[group] (0.85,2.65) rectangle (3.7,3.55);}\visible<4->{\draw (2.2800000000000002,2.15) ellipse (0.73 and .48);\Obj{C}{2.0}{2.05}{1}\Obj{C}{2.56}{2.05}{2}}\visible<2->{\Obj{C}{5.6}{3.0}{1}\Obj{C}{6.14}{3.0}{3}}\visible<2->{\Obj{C}{7.2}{3.0}{3}\Obj{C}{7.74}{3.0}{1}}\visible<3->{\draw[group] (5.250000000000001,2.65) rectangle (8.100000000000001,3.55);}\visible<4->{\draw (6.680000000000001,2.15) ellipse (0.73 and .48);\Obj{C}{6.4}{2.05}{1}\Obj{C}{6.96}{2.05}{3}}\visible<2->{\Obj{C}{10.0}{3.0}{1}\Obj{C}{10.54}{3.0}{4}}\visible<2->{\Obj{C}{11.6}{3.0}{4}\Obj{C}{12.14}{3.0}{1}}\visible<3->{\draw[group] (9.65,2.65) rectangle (12.5,3.55);}\visible<4->{\draw (11.08,2.15) ellipse (0.73 and .48);\Obj{C}{10.8}{2.05}{1}\Obj{C}{11.36}{2.05}{4}}\visible<2->{\Obj{C}{1.2}{1.1}{2}\Obj{C}{1.74}{1.1}{3}}\visible<2->{\Obj{C}{2.8}{1.1}{3}\Obj{C}{3.34}{1.1}{2}}\visible<3->{\draw[group] (0.85,0.7500000000000001) rectangle (3.7,1.6500000000000001);}\visible<4->{\draw (2.2800000000000002,0.2500000000000001) ellipse (0.73 and .48);\Obj{C}{2.0}{0.15000000000000013}{2}\Obj{C}{2.56}{0.15000000000000013}{3}}\visible<2->{\Obj{C}{5.6}{1.1}{2}\Obj{C}{6.14}{1.1}{4}}\visible<2->{\Obj{C}{7.2}{1.1}{4}\Obj{C}{7.74}{1.1}{2}}\visible<3->{\draw[group] (5.250000000000001,0.7500000000000001) rectangle (8.100000000000001,1.6500000000000001);}\visible<4->{\draw (6.680000000000001,0.2500000000000001) ellipse (0.73 and .48);\Obj{C}{6.4}{0.15000000000000013}{2}\Obj{C}{6.96}{0.15000000000000013}{4}}\visible<2->{\Obj{C}{10.0}{1.1}{3}\Obj{C}{10.54}{1.1}{4}}\visible<2->{\Obj{C}{11.6}{1.1}{4}\Obj{C}{12.14}{1.1}{3}}\visible<3->{\draw[group] (9.65,0.7500000000000001) rectangle (12.5,1.6500000000000001);}\visible<4->{\draw (11.08,0.2500000000000001) ellipse (0.73 and .48);\Obj{C}{10.8}{0.15000000000000013}{3}\Obj{C}{11.36}{0.15000000000000013}{4}}
\DiagramEnd
\begin{overlayarea}{\linewidth}{1.15cm}\only<3>{12 ordered lists; two lists in each membership box.\par}
\only<4>{6 groups. An oval has members, with no labelled positions.\par}
\only<5->{$12\div2!=6$ distinct groups.\par}\end{overlayarea}
\end{frame}

% Frame 38. Teaching note: The two 2! factors happen to have the same value but count different things.
\begin{frame}[t]{4B. Write the same division with factorials}
\Q{What does each divisor remove?}
\DiagramStart
\visible<2->{\node[] at (6.75,2.7) {$\displaystyle\frac{4!}{2!}=12\quad\hbox{ordered lists}$};}\visible<3->{\node[] at (6.75,1.2) {$\displaystyle\frac{4!}{2!\,2!}=12\div2=6\quad\hbox{groups}$};}
\DiagramEnd
\begin{overlayarea}{\linewidth}{1.15cm}\only<2>{The unused-tail orders were removed to get the twelve lists.\par}
\only<3->{Now remove the two list orders of each selected pair as well.\par}\end{overlayarea}
\end{frame}

% Frame 39. Teaching note: Collect predictions before revealing outcomes. Numbers identify objects; never multiply the identifiers.
\begin{frame}[t]{4C. Choose three members from four}
\Q{Predict: how many different outcomes are possible?}
\DiagramStart
\Obj{C}{4.8}{2.8}{1}\Obj{C}{5.9}{2.8}{2}\Obj{C}{7.0}{2.8}{3}\Obj{C}{8.1}{2.8}{4}\node[] at (6.7,1.7) {Available distinct objects};\node[] at (6.75,0.5) {One group contains three distinct members, with no roles.};
\DiagramEnd
\begin{overlayarea}{\linewidth}{1.15cm}\end{overlayarea}
\end{frame}

% Frame 40. Teaching note: These are orders of the selected members. None of the six lists makes a different membership.
\begin{frame}[t]{4C. Develop one membership block: 1, 2, 3}
\Q{How many lists can name this one group?}
\DiagramStart
\visible<2->{\Obj{C}{1.4}{2.8}{1}\Obj{C}{1.94}{2.8}{2}\Obj{C}{2.48}{2.8}{3}}\visible<2->{\Obj{C}{1.4}{1.5499999999999998}{1}\Obj{C}{1.94}{1.5499999999999998}{3}\Obj{C}{2.48}{1.5499999999999998}{2}}\visible<3->{\Obj{C}{5.6}{2.8}{2}\Obj{C}{6.14}{2.8}{1}\Obj{C}{6.68}{2.8}{3}}\visible<3->{\Obj{C}{5.6}{1.5499999999999998}{2}\Obj{C}{6.14}{1.5499999999999998}{3}\Obj{C}{6.68}{1.5499999999999998}{1}}\visible<4->{\Obj{C}{9.8}{2.8}{3}\Obj{C}{10.34}{2.8}{1}\Obj{C}{10.88}{2.8}{2}}\visible<4->{\Obj{C}{9.8}{1.5499999999999998}{3}\Obj{C}{10.34}{1.5499999999999998}{2}\Obj{C}{10.88}{1.5499999999999998}{1}}\visible<5->{\draw[group] (1.05,1.05) rectangle (2.9,3.35);}\visible<5->{\draw[group] (5.25,1.05) rectangle (7.1,3.35);}\visible<5->{\draw[group] (9.450000000000001,1.05) rectangle (11.3,3.35);}
\DiagramEnd
\begin{overlayarea}{\linewidth}{1.15cm}\only<5>{Three first-member columns; two orders of the remaining members.\par}
\only<6->{$3\times2\times1=3!=6$ lists, all naming one group.\par}\end{overlayarea}
\end{frame}

% Frame 41. Teaching note: The two frames together show all 24 ordered triples. Ovals are one-per-block summaries.
\begin{frame}[t]{4C. Two membership blocks}
\Q{Does changing the list change the members?}
\DiagramStart
\node[font=\small] at (2.9,3.55) {Members 1, 2, 3};\visible<2->{\Obj{C}{1.0}{2.7}{1}\Obj{C}{1.47}{2.7}{2}\Obj{C}{1.94}{2.7}{3}}\visible<2->{\Obj{C}{1.0}{1.7000000000000002}{1}\Obj{C}{1.47}{1.7000000000000002}{3}\Obj{C}{1.94}{1.7000000000000002}{2}}\visible<2->{\Obj{C}{2.75}{2.7}{2}\Obj{C}{3.22}{2.7}{1}\Obj{C}{3.69}{2.7}{3}}\visible<2->{\Obj{C}{2.75}{1.7000000000000002}{2}\Obj{C}{3.22}{1.7000000000000002}{3}\Obj{C}{3.69}{1.7000000000000002}{1}}\visible<2->{\Obj{C}{4.5}{2.7}{3}\Obj{C}{4.97}{2.7}{1}\Obj{C}{5.44}{2.7}{2}}\visible<2->{\Obj{C}{4.5}{1.7000000000000002}{3}\Obj{C}{4.97}{1.7000000000000002}{2}\Obj{C}{5.44}{1.7000000000000002}{1}}\visible<4->{\draw[group] (0.7,1.1) rectangle (5.85,3.25);}\visible<5->{\draw (3.21,0.44999999999999996) ellipse (1.01 and .48);\Obj{C}{2.65}{0.35}{1}\Obj{C}{3.21}{0.35}{2}\Obj{C}{3.77}{0.35}{3}}\node[font=\small] at (9.4,3.55) {Members 1, 2, 4};\visible<3->{\Obj{C}{7.5}{2.7}{1}\Obj{C}{7.97}{2.7}{2}\Obj{C}{8.44}{2.7}{4}}\visible<3->{\Obj{C}{7.5}{1.7000000000000002}{1}\Obj{C}{7.97}{1.7000000000000002}{4}\Obj{C}{8.44}{1.7000000000000002}{2}}\visible<3->{\Obj{C}{9.25}{2.7}{2}\Obj{C}{9.72}{2.7}{1}\Obj{C}{10.19}{2.7}{4}}\visible<3->{\Obj{C}{9.25}{1.7000000000000002}{2}\Obj{C}{9.72}{1.7000000000000002}{4}\Obj{C}{10.19}{1.7000000000000002}{1}}\visible<3->{\Obj{C}{11.0}{2.7}{4}\Obj{C}{11.47}{2.7}{1}\Obj{C}{11.94}{2.7}{2}}\visible<3->{\Obj{C}{11.0}{1.7000000000000002}{4}\Obj{C}{11.47}{1.7000000000000002}{2}\Obj{C}{11.94}{1.7000000000000002}{1}}\visible<4->{\draw[group] (7.2,1.1) rectangle (12.35,3.25);}\visible<5->{\draw (9.71,0.44999999999999996) ellipse (1.01 and .48);\Obj{C}{9.15}{0.35}{1}\Obj{C}{9.71}{0.35}{2}\Obj{C}{10.27}{0.35}{4}}
\DiagramEnd
\begin{overlayarea}{\linewidth}{1.15cm}\only<4>{Each block contains six lists of its own three members.\par}
\only<5->{Each block collapses to one membership diagram.\par}\end{overlayarea}
\end{frame}

% Frame 42. Teaching note: The two frames together show all 24 ordered triples. Ovals are one-per-block summaries.
\begin{frame}[t]{4C. Two membership blocks}
\Q{Does changing the list change the members?}
\DiagramStart
\node[font=\small] at (2.9,3.55) {Members 1, 3, 4};\visible<2->{\Obj{C}{1.0}{2.7}{1}\Obj{C}{1.47}{2.7}{3}\Obj{C}{1.94}{2.7}{4}}\visible<2->{\Obj{C}{1.0}{1.7000000000000002}{1}\Obj{C}{1.47}{1.7000000000000002}{4}\Obj{C}{1.94}{1.7000000000000002}{3}}\visible<2->{\Obj{C}{2.75}{2.7}{3}\Obj{C}{3.22}{2.7}{1}\Obj{C}{3.69}{2.7}{4}}\visible<2->{\Obj{C}{2.75}{1.7000000000000002}{3}\Obj{C}{3.22}{1.7000000000000002}{4}\Obj{C}{3.69}{1.7000000000000002}{1}}\visible<2->{\Obj{C}{4.5}{2.7}{4}\Obj{C}{4.97}{2.7}{1}\Obj{C}{5.44}{2.7}{3}}\visible<2->{\Obj{C}{4.5}{1.7000000000000002}{4}\Obj{C}{4.97}{1.7000000000000002}{3}\Obj{C}{5.44}{1.7000000000000002}{1}}\visible<4->{\draw[group] (0.7,1.1) rectangle (5.85,3.25);}\visible<5->{\draw (3.21,0.44999999999999996) ellipse (1.01 and .48);\Obj{C}{2.65}{0.35}{1}\Obj{C}{3.21}{0.35}{3}\Obj{C}{3.77}{0.35}{4}}\node[font=\small] at (9.4,3.55) {Members 2, 3, 4};\visible<3->{\Obj{C}{7.5}{2.7}{2}\Obj{C}{7.97}{2.7}{3}\Obj{C}{8.44}{2.7}{4}}\visible<3->{\Obj{C}{7.5}{1.7000000000000002}{2}\Obj{C}{7.97}{1.7000000000000002}{4}\Obj{C}{8.44}{1.7000000000000002}{3}}\visible<3->{\Obj{C}{9.25}{2.7}{3}\Obj{C}{9.72}{2.7}{2}\Obj{C}{10.19}{2.7}{4}}\visible<3->{\Obj{C}{9.25}{1.7000000000000002}{3}\Obj{C}{9.72}{1.7000000000000002}{4}\Obj{C}{10.19}{1.7000000000000002}{2}}\visible<3->{\Obj{C}{11.0}{2.7}{4}\Obj{C}{11.47}{2.7}{2}\Obj{C}{11.94}{2.7}{3}}\visible<3->{\Obj{C}{11.0}{1.7000000000000002}{4}\Obj{C}{11.47}{1.7000000000000002}{3}\Obj{C}{11.94}{1.7000000000000002}{2}}\visible<4->{\draw[group] (7.2,1.1) rectangle (12.35,3.25);}\visible<5->{\draw (9.71,0.44999999999999996) ellipse (1.01 and .48);\Obj{C}{9.15}{0.35}{2}\Obj{C}{9.71}{0.35}{3}\Obj{C}{10.27}{0.35}{4}}
\DiagramEnd
\begin{overlayarea}{\linewidth}{1.15cm}\only<4>{Each block contains six lists of its own three members.\par}
\only<5->{Each block collapses to one membership diagram.\par}\end{overlayarea}
\end{frame}

% Frame 43. Teaching note: Division works because every block has equal size, not because lists are simply discarded.
\begin{frame}[t]{4C. Four blocks, six lists in each}
\Q{How do twenty-four lists become four groups?}
\DiagramStart
\visible<2->{\draw[group] (0.5,1.1) rectangle (2.8,3.4);}\visible<2->{\draw (1.6600000000000001,2.7) ellipse (1.01 and .48);\Obj{C}{1.1}{2.6}{1}\Obj{C}{1.66}{2.6}{2}\Obj{C}{2.22}{2.6}{3}}\visible<2->{\node[font=\small] at (1.65,1.5) {6 lists};}\visible<2->{\draw[group] (3.8500000000000005,1.1) rectangle (6.15,3.4);}\visible<2->{\draw (5.01,2.7) ellipse (1.01 and .48);\Obj{C}{4.45}{2.6}{1}\Obj{C}{5.01}{2.6}{2}\Obj{C}{5.57}{2.6}{4}}\visible<2->{\node[font=\small] at (5.0,1.5) {6 lists};}\visible<2->{\draw[group] (7.2,1.1) rectangle (9.5,3.4);}\visible<2->{\draw (8.36,2.7) ellipse (1.01 and .48);\Obj{C}{7.8}{2.6}{1}\Obj{C}{8.36}{2.6}{3}\Obj{C}{8.92}{2.6}{4}}\visible<2->{\node[font=\small] at (8.35,1.5) {6 lists};}\visible<2->{\draw[group] (10.55,1.1) rectangle (12.850000000000001,3.4);}\visible<2->{\draw (11.710000000000003,2.7) ellipse (1.01 and .48);\Obj{C}{11.15}{2.6}{2}\Obj{C}{11.71}{2.6}{3}\Obj{C}{12.27}{2.6}{4}}\visible<2->{\node[font=\small] at (11.700000000000001,1.5) {6 lists};}
\DiagramEnd
\begin{overlayarea}{\linewidth}{1.15cm}\only<2>{Four membership blocks have the same six-list structure.\par}
\only<3->{$24$ ordered lists $\div6$ lists per group $=4$ groups.\par}\end{overlayarea}
\end{frame}

% Frame 44. Teaching note: n counts available distinct objects; r counts members, 0<=r<=n. The denominator factors remove different overcounts.
\begin{frame}[t]{Groups: divide by orders of the selected members}
\Q{How many lists name every group of r members?}
\DiagramStart
\visible<2->{\node[] at (6.75,2.9) {Every $r$-member group has $r!$ lists.};}\visible<3->{\node[] at (6.75,1.4) {$\displaystyle\binom nr=\frac{\text{ordered choices}}{\text{orders per group}}=\frac{n!}{r!(n-r)!}$};}\visible<4->{\node[] at (6.75,0.2) {Read the symbol as ``$n$ choose $r$.};}
\DiagramEnd
\begin{overlayarea}{\linewidth}{1.15cm}\only<3>{$(n-r)!$ removes orders of the unused objects.\par}
\only<4->{$r!$ removes list orders of the selected members. Membership stays the same.\par}\end{overlayarea}
\end{frame}

\section{5. Arrange indistinguishable copies}
% Frame 45. Teaching note: Copies begin without numbers; temporary identifiers appear only in the next labelled counting stage.
\begin{frame}[t]{5A. Two matching circles and one triangle}
\Q{Predict: how many different visible rows can we make?}
\DiagramStart
\Obj{C}{5.0}{2.6}{}\Obj{C}{6.1}{2.6}{}\Obj{T}{7.2}{2.6}{}\node[] at (6.4,1.2) {Place every supplied copy exactly once.};
\DiagramEnd
\begin{overlayarea}{\linewidth}{1.15cm}\only<2->{Positions matter. Exchanging matching circles makes no visible change.\par}\end{overlayarea}
\end{frame}

% Frame 46. Teaching note: Overlay 5 removes labels without moving shapes. Overlay 6 encloses equal pairs after the change is visible.
\begin{frame}[t]{5A. Count with temporary circle identifiers}
\Q{Which labelled rows become identical when identifiers vanish?}
\DiagramStart
\node[font=\small] at (1.95,3.55) {Triangle last};\visible<2->{\Obj{C}{1.4}{2.7}{\alt<5->{}{1}}\Obj{C}{1.95}{2.7}{\alt<5->{}{2}}\Obj{T}{2.5}{2.7}{}}\visible<2->{\Obj{C}{1.4}{1.4500000000000002}{\alt<5->{}{2}}\Obj{C}{1.95}{1.4500000000000002}{\alt<5->{}{1}}\Obj{T}{2.5}{1.4500000000000002}{}}\visible<6->{\draw[group] (1.0499999999999998,0.95) rectangle (2.8499999999999996,3.25);}\node[font=\small] at (6.1499999999999995,3.55) {Triangle middle};\visible<3->{\Obj{C}{5.6}{2.7}{\alt<5->{}{1}}\Obj{T}{6.1499999999999995}{2.7}{}\Obj{C}{6.699999999999999}{2.7}{\alt<5->{}{2}}}\visible<3->{\Obj{C}{5.6}{1.4500000000000002}{\alt<5->{}{2}}\Obj{T}{6.1499999999999995}{1.4500000000000002}{}\Obj{C}{6.699999999999999}{1.4500000000000002}{\alt<5->{}{1}}}\visible<6->{\draw[group] (5.25,0.95) rectangle (7.05,3.25);}\node[font=\small] at (10.350000000000001,3.55) {Triangle first};\visible<4->{\Obj{T}{9.8}{2.7}{}\Obj{C}{10.350000000000001}{2.7}{\alt<5->{}{1}}\Obj{C}{10.9}{2.7}{\alt<5->{}{2}}}\visible<4->{\Obj{T}{9.8}{1.4500000000000002}{}\Obj{C}{10.350000000000001}{1.4500000000000002}{\alt<5->{}{2}}\Obj{C}{10.9}{1.4500000000000002}{\alt<5->{}{1}}}\visible<6->{\draw[group] (9.450000000000001,0.95) rectangle (11.25,3.25);}
\DiagramEnd
\begin{overlayarea}{\linewidth}{1.15cm}\only<4>{Three labelled objects give $3!=6$ rows.\par}
\only<5>{Erase only the temporary identifiers. Both rows remain visible.\par}
\only<6->{Each outlined pair is now two copies of the same visible row.\par}\end{overlayarea}
\end{frame}

% Frame 47. Teaching note: This new frame collapses the three stationary pairs to one row each; identifiers remain absent.
\begin{frame}[t]{5A. Keep one row from each matching pair}
\Q{Why divide by two circle-label orders?}
\DiagramStart
\Obj{C}{1.4}{2.5}{}\Obj{C}{1.94}{2.5}{}\Obj{T}{2.48}{2.5}{}\Obj{C}{5.6}{2.5}{}\Obj{T}{6.14}{2.5}{}\Obj{C}{6.68}{2.5}{}\Obj{T}{9.8}{2.5}{}\Obj{C}{10.34}{2.5}{}\Obj{C}{10.88}{2.5}{}
\DiagramEnd
\begin{overlayarea}{\linewidth}{1.15cm}\only<2>{Every visible row had $2!=2$ labelled versions.\par}
\only<3->{$3!\div2!=6\div2=3$ visible arrangements.\par}\end{overlayarea}
\end{frame}

% Frame 48. Teaching note: Every supplied object is placed once; circle positions determine the visible arrangement.
\begin{frame}[t]{5B. Two matching circles and two matching triangles}
\Q{Predict: how many visible rows use these four copies?}
\DiagramStart
\Obj{C}{4.5}{2.5}{}\Obj{C}{5.6}{2.5}{}\Obj{T}{6.7}{2.5}{}\Obj{T}{7.8}{2.5}{}
\DiagramEnd
\begin{overlayarea}{\linewidth}{1.15cm}\end{overlayarea}
\end{frame}

% Frame 49. Teaching note: Position numbers below the rows identify slots. There are no numbers above the real matching copies.
\begin{frame}[t]{5B. Organize by the two circle positions}
\Q{How many labelled versions belong to each visible row?}
\DiagramStart
\visible<2->{\Obj{C}{1.1}{3.0}{}\node[font=\small] at (1.1,2.68) {1};\Obj{C}{1.65}{3.0}{}\node[font=\small] at (1.65,2.68) {2};\Obj{T}{2.2}{3.0}{}\node[font=\small] at (2.2,2.68) {3};\Obj{T}{2.75}{3.0}{}\node[font=\small] at (2.75,2.68) {4};}\visible<2->{\node[font=\small] at (1.9000000000000001,2.27) {Circle positions 1, 2};}\visible<2->{\Obj{C}{5.5}{3.0}{}\node[font=\small] at (5.5,2.68) {1};\Obj{T}{6.05}{3.0}{}\node[font=\small] at (6.05,2.68) {2};\Obj{C}{6.6}{3.0}{}\node[font=\small] at (6.6,2.68) {3};\Obj{T}{7.15}{3.0}{}\node[font=\small] at (7.15,2.68) {4};}\visible<2->{\node[font=\small] at (6.3,2.27) {Circle positions 1, 3};}\visible<2->{\Obj{C}{9.9}{3.0}{}\node[font=\small] at (9.9,2.68) {1};\Obj{T}{10.45}{3.0}{}\node[font=\small] at (10.45,2.68) {2};\Obj{T}{11.0}{3.0}{}\node[font=\small] at (11.0,2.68) {3};\Obj{C}{11.55}{3.0}{}\node[font=\small] at (11.55,2.68) {4};}\visible<2->{\node[font=\small] at (10.700000000000001,2.27) {Circle positions 1, 4};}\visible<2->{\Obj{T}{1.1}{1.1}{}\node[font=\small] at (1.1,0.78) {1};\Obj{C}{1.65}{1.1}{}\node[font=\small] at (1.65,0.78) {2};\Obj{C}{2.2}{1.1}{}\node[font=\small] at (2.2,0.78) {3};\Obj{T}{2.75}{1.1}{}\node[font=\small] at (2.75,0.78) {4};}\visible<2->{\node[font=\small] at (1.9000000000000001,0.3700000000000001) {Circle positions 2, 3};}\visible<2->{\Obj{T}{5.5}{1.1}{}\node[font=\small] at (5.5,0.78) {1};\Obj{C}{6.05}{1.1}{}\node[font=\small] at (6.05,0.78) {2};\Obj{T}{6.6}{1.1}{}\node[font=\small] at (6.6,0.78) {3};\Obj{C}{7.15}{1.1}{}\node[font=\small] at (7.15,0.78) {4};}\visible<2->{\node[font=\small] at (6.3,0.3700000000000001) {Circle positions 2, 4};}\visible<2->{\Obj{T}{9.9}{1.1}{}\node[font=\small] at (9.9,0.78) {1};\Obj{T}{10.45}{1.1}{}\node[font=\small] at (10.45,0.78) {2};\Obj{C}{11.0}{1.1}{}\node[font=\small] at (11.0,0.78) {3};\Obj{C}{11.55}{1.1}{}\node[font=\small] at (11.55,0.78) {4};}\visible<2->{\node[font=\small] at (10.700000000000001,0.3700000000000001) {Circle positions 3, 4};}
\DiagramEnd
\begin{overlayarea}{\linewidth}{1.15cm}\only<3->{These six rows use each possible pair of circle positions once.\par}\end{overlayarea}
\end{frame}

% Frame 50. Teaching note: Row changes are independent of column changes. Temporary identifiers disappear at overlay 5; shapes never move.
\begin{frame}[t]{5B. Expand one visible row into labelled versions}
\Q{What changes across rows? What changes across columns?}
\DiagramStart
\Obj{C}{2.0}{2.8}{\alt<5->{}{1}}\Obj{C}{2.75}{2.8}{\alt<5->{}{2}}\Obj{T}{3.5}{2.8}{\alt<5->{}{1}}\Obj{T}{4.25}{2.8}{\alt<5->{}{2}}\visible<5->{\draw[group] (1.65,2.4) rectangle (4.6,3.4);}\Obj{C}{8.2}{2.8}{\alt<5->{}{1}}\Obj{C}{8.95}{2.8}{\alt<5->{}{2}}\Obj{T}{9.7}{2.8}{\alt<5->{}{2}}\Obj{T}{10.45}{2.8}{\alt<5->{}{1}}\visible<5->{\draw[group] (7.85,2.4) rectangle (10.799999999999999,3.4);}\Obj{C}{2.0}{1.15}{\alt<5->{}{2}}\Obj{C}{2.75}{1.15}{\alt<5->{}{1}}\Obj{T}{3.5}{1.15}{\alt<5->{}{1}}\Obj{T}{4.25}{1.15}{\alt<5->{}{2}}\visible<5->{\draw[group] (1.65,0.7499999999999999) rectangle (4.6,1.75);}\Obj{C}{8.2}{1.15}{\alt<5->{}{2}}\Obj{C}{8.95}{1.15}{\alt<5->{}{1}}\Obj{T}{9.7}{1.15}{\alt<5->{}{2}}\Obj{T}{10.45}{1.15}{\alt<5->{}{1}}\visible<5->{\draw[group] (7.85,0.7499999999999999) rectangle (10.799999999999999,1.75);}
\DiagramEnd
\begin{overlayarea}{\linewidth}{1.15cm}\only<2>{Rows switch circle identifiers; columns switch triangle identifiers.\par}
\only<3-4>{$2!$ circle orders $\times2!$ triangle orders $=4$ labelled versions.\par}
\only<5->{Erase identifiers: all four entries become the same visible row.\par}\end{overlayarea}
\end{frame}

% Frame 51. Teaching note: Row changes are independent of column changes. Temporary identifiers disappear at overlay 5; shapes never move.
\begin{frame}[t]{5B. Expand one visible row into labelled versions}
\Q{What changes across rows? What changes across columns?}
\DiagramStart
\Obj{C}{2.0}{2.8}{\alt<5->{}{1}}\Obj{T}{2.75}{2.8}{\alt<5->{}{1}}\Obj{C}{3.5}{2.8}{\alt<5->{}{2}}\Obj{T}{4.25}{2.8}{\alt<5->{}{2}}\visible<5->{\draw[group] (1.65,2.4) rectangle (4.6,3.4);}\Obj{C}{8.2}{2.8}{\alt<5->{}{1}}\Obj{T}{8.95}{2.8}{\alt<5->{}{2}}\Obj{C}{9.7}{2.8}{\alt<5->{}{2}}\Obj{T}{10.45}{2.8}{\alt<5->{}{1}}\visible<5->{\draw[group] (7.85,2.4) rectangle (10.799999999999999,3.4);}\Obj{C}{2.0}{1.15}{\alt<5->{}{2}}\Obj{T}{2.75}{1.15}{\alt<5->{}{1}}\Obj{C}{3.5}{1.15}{\alt<5->{}{1}}\Obj{T}{4.25}{1.15}{\alt<5->{}{2}}\visible<5->{\draw[group] (1.65,0.7499999999999999) rectangle (4.6,1.75);}\Obj{C}{8.2}{1.15}{\alt<5->{}{2}}\Obj{T}{8.95}{1.15}{\alt<5->{}{2}}\Obj{C}{9.7}{1.15}{\alt<5->{}{1}}\Obj{T}{10.45}{1.15}{\alt<5->{}{1}}\visible<5->{\draw[group] (7.85,0.7499999999999999) rectangle (10.799999999999999,1.75);}
\DiagramEnd
\begin{overlayarea}{\linewidth}{1.15cm}\only<2>{Rows switch circle identifiers; columns switch triangle identifiers.\par}
\only<3-4>{$2!$ circle orders $\times2!$ triangle orders $=4$ labelled versions.\par}
\only<5->{Erase identifiers: all four entries become the same visible row.\par}\end{overlayarea}
\end{frame}

% Frame 52. Teaching note: Each of the six boxes refers to four labelled versions, not four distinct visible arrangements.
\begin{frame}[t]{5B. Every visible row has the same four-version structure}
\Q{Can we use the same divisor for all six rows?}
\DiagramStart
\visible<2->{\Obj{C}{1.1}{3.0}{}\node[font=\small] at (1.1,2.68) {1};\Obj{C}{1.65}{3.0}{}\node[font=\small] at (1.65,2.68) {2};\Obj{T}{2.2}{3.0}{}\node[font=\small] at (2.2,2.68) {3};\Obj{T}{2.75}{3.0}{}\node[font=\small] at (2.75,2.68) {4};}\visible<3->{\draw[group] (0.7500000000000001,2.35) rectangle (3.1,3.5);}\visible<3->{\node[font=\small] at (1.9000000000000001,1.97) {4 labelled versions};}\visible<2->{\Obj{C}{5.5}{3.0}{}\node[font=\small] at (5.5,2.68) {1};\Obj{T}{6.05}{3.0}{}\node[font=\small] at (6.05,2.68) {2};\Obj{C}{6.6}{3.0}{}\node[font=\small] at (6.6,2.68) {3};\Obj{T}{7.15}{3.0}{}\node[font=\small] at (7.15,2.68) {4};}\visible<3->{\draw[group] (5.15,2.35) rectangle (7.5,3.5);}\visible<3->{\node[font=\small] at (6.3,1.97) {4 labelled versions};}\visible<2->{\Obj{C}{9.9}{3.0}{}\node[font=\small] at (9.9,2.68) {1};\Obj{T}{10.45}{3.0}{}\node[font=\small] at (10.45,2.68) {2};\Obj{T}{11.0}{3.0}{}\node[font=\small] at (11.0,2.68) {3};\Obj{C}{11.55}{3.0}{}\node[font=\small] at (11.55,2.68) {4};}\visible<3->{\draw[group] (9.55,2.35) rectangle (11.9,3.5);}\visible<3->{\node[font=\small] at (10.700000000000001,1.97) {4 labelled versions};}\visible<2->{\Obj{T}{1.1}{1.1}{}\node[font=\small] at (1.1,0.78) {1};\Obj{C}{1.65}{1.1}{}\node[font=\small] at (1.65,0.78) {2};\Obj{C}{2.2}{1.1}{}\node[font=\small] at (2.2,0.78) {3};\Obj{T}{2.75}{1.1}{}\node[font=\small] at (2.75,0.78) {4};}\visible<3->{\draw[group] (0.7500000000000001,0.45000000000000007) rectangle (3.1,1.6);}\visible<3->{\node[font=\small] at (1.9000000000000001,0.07000000000000006) {4 labelled versions};}\visible<2->{\Obj{T}{5.5}{1.1}{}\node[font=\small] at (5.5,0.78) {1};\Obj{C}{6.05}{1.1}{}\node[font=\small] at (6.05,0.78) {2};\Obj{T}{6.6}{1.1}{}\node[font=\small] at (6.6,0.78) {3};\Obj{C}{7.15}{1.1}{}\node[font=\small] at (7.15,0.78) {4};}\visible<3->{\draw[group] (5.15,0.45000000000000007) rectangle (7.5,1.6);}\visible<3->{\node[font=\small] at (6.3,0.07000000000000006) {4 labelled versions};}\visible<2->{\Obj{T}{9.9}{1.1}{}\node[font=\small] at (9.9,0.78) {1};\Obj{T}{10.45}{1.1}{}\node[font=\small] at (10.45,0.78) {2};\Obj{C}{11.0}{1.1}{}\node[font=\small] at (11.0,0.78) {3};\Obj{C}{11.55}{1.1}{}\node[font=\small] at (11.55,0.78) {4};}\visible<3->{\draw[group] (9.55,0.45000000000000007) rectangle (11.9,1.6);}\visible<3->{\node[font=\small] at (10.700000000000001,0.07000000000000006) {4 labelled versions};}
\DiagramEnd
\begin{overlayarea}{\linewidth}{1.15cm}\only<3>{Yes. In each box: two circle orders and two triangle orders.\par}
\only<4->{$4!\div(2!\times2!)=24\div4=6$ visible arrangements.\par}\end{overlayarea}
\end{frame}

% Frame 53. Teaching note: Triangle position will organize the complete list; temporary identifiers are not yet present.
\begin{frame}[t]{5C. Three matching circles and one triangle}
\Q{Predict: how many visible rows use these four copies?}
\DiagramStart
\Obj{C}{4.5}{2.6}{}\Obj{C}{5.6}{2.6}{}\Obj{C}{6.7}{2.6}{}\Obj{T}{7.8}{2.6}{}
\DiagramEnd
\begin{overlayarea}{\linewidth}{1.15cm}\end{overlayarea}
\end{frame}

% Frame 54. Teaching note: All four visible rows are complete. Circle copies have no distinguishing numbers.
\begin{frame}[t]{5C. Four triangle positions}
\Q{How many labelled versions hide behind each visible row?}
\DiagramStart
\visible<2->{\Obj{T}{0.8}{2.5}{}\node[font=\small] at (0.8,2.18) {1};\Obj{C}{1.35}{2.5}{}\node[font=\small] at (1.35,2.18) {2};\Obj{C}{1.9}{2.5}{}\node[font=\small] at (1.9,2.18) {3};\Obj{C}{2.45}{2.5}{}\node[font=\small] at (2.45,2.18) {4};}\visible<2->{\node[font=\small,text width=2.6cm,align=center] at (1.6,1.5) {Triangle: position 1};}\visible<3->{\Obj{C}{4.15}{2.5}{}\node[font=\small] at (4.15,2.18) {1};\Obj{T}{4.7}{2.5}{}\node[font=\small] at (4.7,2.18) {2};\Obj{C}{5.25}{2.5}{}\node[font=\small] at (5.25,2.18) {3};\Obj{C}{5.8}{2.5}{}\node[font=\small] at (5.8,2.18) {4};}\visible<3->{\node[font=\small,text width=2.6cm,align=center] at (4.95,1.5) {Triangle: position 2};}\visible<4->{\Obj{C}{7.5}{2.5}{}\node[font=\small] at (7.5,2.18) {1};\Obj{C}{8.05}{2.5}{}\node[font=\small] at (8.05,2.18) {2};\Obj{T}{8.6}{2.5}{}\node[font=\small] at (8.6,2.18) {3};\Obj{C}{9.15}{2.5}{}\node[font=\small] at (9.15,2.18) {4};}\visible<4->{\node[font=\small,text width=2.6cm,align=center] at (8.3,1.5) {Triangle: position 3};}\visible<5->{\Obj{C}{10.85}{2.5}{}\node[font=\small] at (10.85,2.18) {1};\Obj{C}{11.4}{2.5}{}\node[font=\small] at (11.4,2.18) {2};\Obj{C}{11.95}{2.5}{}\node[font=\small] at (11.95,2.18) {3};\Obj{T}{12.5}{2.5}{}\node[font=\small] at (12.5,2.18) {4};}\visible<5->{\node[font=\small,text width=2.6cm,align=center] at (11.650000000000002,1.5) {Triangle: position 4};}
\DiagramEnd
\begin{overlayarea}{\linewidth}{1.15cm}\only<6->{Four visible arrangements. Keep the triangle fixed and label only the circles.\par}\end{overlayarea}
\end{frame}

% Frame 55. Teaching note: All six label orders are shown in three columns of two. Triangle stays in the fourth position.
\begin{frame}[t]{5C. Fix the triangle last; order the circle labels}
\Q{How many circle-label orders give this one visible row?}
\DiagramStart
\visible<2->{\Obj{C}{1.0}{2.8}{1}\Obj{C}{1.6}{2.8}{2}\Obj{C}{2.2}{2.8}{3}\Obj{T}{2.8}{2.8}{}}\visible<2->{\Obj{C}{1.0}{1.4499999999999997}{1}\Obj{C}{1.6}{1.4499999999999997}{3}\Obj{C}{2.2}{1.4499999999999997}{2}\Obj{T}{2.8}{1.4499999999999997}{}}\visible<3->{\Obj{C}{5.3}{2.8}{2}\Obj{C}{5.9}{2.8}{1}\Obj{C}{6.5}{2.8}{3}\Obj{T}{7.1}{2.8}{}}\visible<3->{\Obj{C}{5.3}{1.4499999999999997}{2}\Obj{C}{5.9}{1.4499999999999997}{3}\Obj{C}{6.5}{1.4499999999999997}{1}\Obj{T}{7.1}{1.4499999999999997}{}}\visible<4->{\Obj{C}{9.6}{2.8}{3}\Obj{C}{10.2}{2.8}{1}\Obj{C}{10.8}{2.8}{2}\Obj{T}{11.4}{2.8}{}}\visible<4->{\Obj{C}{9.6}{1.4499999999999997}{3}\Obj{C}{10.2}{1.4499999999999997}{2}\Obj{C}{10.8}{1.4499999999999997}{1}\Obj{T}{11.4}{1.4499999999999997}{}}\visible<5->{\draw[group] (0.65,0.95) rectangle (3.1,3.4);}\visible<5->{\draw[group] (4.95,0.95) rectangle (7.4,3.4);}\visible<5->{\draw[group] (9.25,0.95) rectangle (11.7,3.4);}
\DiagramEnd
\begin{overlayarea}{\linewidth}{1.15cm}\only<5>{Three first-circle groups; two label orders in each.\par}
\only<6->{$3\times2\times1=3!=6$ labelled versions of this visible row.\par}\end{overlayarea}
\end{frame}

% Frame 56. Teaching note: Triangle position does not restrict the independent circle-label orders.
\begin{frame}[t]{5C. The same six-version structure in every box}
\Q{How many labelled rows altogether?}
\DiagramStart
\Obj{T}{0.8}{2.7}{}\Obj{C}{1.35}{2.7}{}\Obj{C}{1.9}{2.7}{}\Obj{C}{2.45}{2.7}{}\visible<2->{\draw[group] (0.5,1) rectangle (2.8,3.25);}\visible<2->{\node[font=\small] at (1.6,1.5) {6 labelled versions};}\Obj{C}{4.15}{2.7}{}\Obj{T}{4.7}{2.7}{}\Obj{C}{5.25}{2.7}{}\Obj{C}{5.8}{2.7}{}\visible<2->{\draw[group] (3.8500000000000005,1) rectangle (6.15,3.25);}\visible<2->{\node[font=\small] at (4.95,1.5) {6 labelled versions};}\Obj{C}{7.5}{2.7}{}\Obj{C}{8.05}{2.7}{}\Obj{T}{8.6}{2.7}{}\Obj{C}{9.15}{2.7}{}\visible<2->{\draw[group] (7.2,1) rectangle (9.5,3.25);}\visible<2->{\node[font=\small] at (8.3,1.5) {6 labelled versions};}\Obj{C}{10.85}{2.7}{}\Obj{C}{11.4}{2.7}{}\Obj{C}{11.95}{2.7}{}\Obj{T}{12.5}{2.7}{}\visible<2->{\draw[group] (10.55,1) rectangle (12.850000000000001,3.25);}\visible<2->{\node[font=\small] at (11.650000000000002,1.5) {6 labelled versions};}
\DiagramEnd
\begin{overlayarea}{\linewidth}{1.15cm}\only<2>{Four visible arrangements; six versions each: $4\times6=24$.\par}
\only<3->{$4!\div3!=24\div6=4$ visible arrangements.\par}\end{overlayarea}
\end{frame}

% Frame 57. Teaching note: Subscripts are identifiers, not multiplication. k is positive and i in m_i ranges from 1 through k.
\begin{frame}[t]{Matching copies: name the copy counts}
\Q{What do the subscripts tell us?}
\DiagramStart
\visible<2->{\node[] at (6.75,2.8) {$n$: total supplied objects; $k$: number of types};}\visible<3->{\node[] at (6.75,1.6) {$m_1,\ m_2,\ \ldots,\ m_k$: copy counts for types $1,2,\ldots,k$};}\visible<4->{\node[] at (6.75,0.4) {$m_1+m_2+\cdots+m_k=n$};}
\DiagramEnd
\begin{overlayarea}{\linewidth}{1.15cm}\only<3>{A subscript identifies a type. For example, the 2 in $m_2$ means type 2.\par}
\only<4->{Every copy count is a positive whole number. The counts add to the total.\par}\end{overlayarea}
\end{frame}

% Frame 58. Teaching note: The dots continue through all k types. Order of positions still matters; only invisible label exchanges are removed.
\begin{frame}[t]{Matching copies: the general rule}
\Q{How often does each visible arrangement get counted?}
\DiagramStart
\visible<2->{\node[] at (6.75,2.9) {$m_1!\times m_2!\times\cdots\times m_k!$ label orders};}\visible<3->{\node[] at (6.75,1.3) {$\displaystyle\text{visible rows}=\frac{n!}{m_1!m_2!\cdots m_k!}$};}
\DiagramEnd
\begin{overlayarea}{\linewidth}{1.15cm}\only<2>{For every label order of one type, all label orders of the other types remain.\par}
\only<3>{So the label orders multiply across types; divide them out of the full count.\par}
\only<4->{Identical supplied copies are different from reuse: place each supplied copy once.\par}\end{overlayarea}
\end{frame}

\section{Synthesis}
% Frame 59. Teaching note: Ask the outcome question first. Membership has no jobs; arrangements keep positions.
\begin{frame}[t]{A visual decision guide}
\Q{Are we choosing from inputs, assigning positions, or choosing members?}
\DiagramStart
\node[] at (6.75,3.5) {What makes an outcome different?};\node[font=\bfseries] at (2.1,2.55) {Input groups};\node[font=\small] at (2.1,1.8) {One choice from each};\node[font=\bfseries] at (6.75,2.55) {Membership};\node[font=\small] at (6.75,1.8) {A group without roles};\node[font=\bfseries] at (11.4,2.55) {Positions};\node[font=\small] at (11.4,1.8) {A row or job assignment};\visible<2->{\node[font=\small] at (2.1,0.6) {Multiply step counts};}\visible<3->{\node[font=\small,text width=3.6cm,align=center] at (6.75,0.6) {Distinct members: divide list orders};}\visible<4->{\node[font=\small,text width=3.3cm,align=center] at (11.4,0.6) {All or some? Matching copies?};}
\DiagramEnd
\begin{overlayarea}{\linewidth}{1.15cm}\only<4->{For positions: all distinct objects, some distinct objects, or all matching copies?\par}\end{overlayarea}
\end{frame}

% Frame 60. Teaching note: The miniature diagrams anchor rules already derived; no new contexts or formulas.
\begin{frame}[t]{Rule summary: what the factors count}
\Q{Can you explain each factor or divisor using its picture?}
\DiagramStart
\node[anchor=west,font=\small] at (0.1,3.2) {Input choices};\Obj{C}{3.1}{3.2}{1}\Obj{T}{3.64}{3.2}{1}\node[] at (6.5,3.2) {$2\times3$};\node[text width=4cm,align=left,font=\small] at (9.8,3.2) {2 circle groups; 3 triangle choices each};\node[anchor=west,font=\small] at (0.1,1.9200000000000002) {All distinct};\Obj{C}{3.1}{1.9200000000000002}{1}\Obj{C}{3.64}{1.9200000000000002}{2}\Obj{C}{4.18}{1.9200000000000002}{3}\node[] at (6.5,1.9200000000000002) {$n!$};\node[text width=4cm,align=left,font=\small] at (9.8,1.9200000000000002) {one fewer choice after every placement};\node[anchor=west,font=\small] at (0.1,0.6400000000000001) {Some ordered};\Obj{C}{3.1}{0.6400000000000001}{1}\Obj{C}{3.64}{0.6400000000000001}{2}\node[] at (6.5,0.6400000000000001) {$n!/(n-r)!$};\node[text width=4cm,align=left,font=\small] at (9.8,0.6400000000000001) {divide by unused-tail orders};
\DiagramEnd
\begin{overlayarea}{\linewidth}{1.15cm}\only<2->{Input counts stay the same after earlier choices. Distinct-object counts decrease.\par}\end{overlayarea}
\end{frame}

% Frame 61. Teaching note: Do not confuse member-list orders, unused-tail orders and temporary-copy-label orders.
\begin{frame}[t]{Rule summary: what division removes}
\Q{Which versions are different lists but the same actual outcome?}
\DiagramStart
\visible<1->{\draw (2.08,2.9) ellipse (0.73 and .48);\Obj{C}{1.8}{2.8}{1}\Obj{C}{2.36}{2.8}{2}}\node[] at (7.9,2.8) {$\displaystyle\binom nr=\frac{n!}{r!(n-r)!}$};\node[font=\small] at (7.9,1.85) {$r!$: selected-member list orders; $(n-r)!$: unused-tail orders};\Obj{C}{1.8}{0.7}{}\Obj{C}{2.34}{0.7}{}\Obj{T}{2.88}{0.7}{}\node[] at (8.2,0.7) {$\displaystyle\frac{n!}{m_1!m_2!\cdots m_k!}$};
\DiagramEnd
\begin{overlayarea}{\linewidth}{1.15cm}\only<2>{Groups: remove list orders of the same members.\par}
\only<3->{Matching copies: remove label orders within each type. Positions still matter.\par}\end{overlayarea}
\end{frame}

\section{Oral checks}
% Frame 62. Teaching note: Wait for reasoning before overlay 2. Identifiers do not represent numeric values.
\begin{frame}[t]{Oral check 1: input groups}
\Q{Two circles and two triangles: how many circle--triangle outcomes?}
\DiagramStart
\Obj{C}{2}{2.7}{1}\Obj{C}{3}{2.7}{2}\Obj{T}{8}{2.7}{1}\Obj{T}{9}{2.7}{2}
\DiagramEnd
\begin{overlayarea}{\linewidth}{1.15cm}\only<2->{$2\times2=4$. Each circle leaves two triangle choices.\par}\end{overlayarea}
\end{frame}

% Frame 63. Teaching note: Do not reveal the member answer together with the job answer.
\begin{frame}[t]{Oral check 2: jobs or members?}
\Q{From objects 1, 2, 3: how many two-job assignments? Two-member groups?}
\DiagramStart
\Obj{C}{5.0}{2.6}{1}\Obj{C}{6.1}{2.6}{2}\Obj{C}{7.2}{2.6}{3}
\DiagramEnd
\begin{overlayarea}{\linewidth}{1.15cm}\only<2>{Jobs: $3\times2=6$. Switching holders changes their jobs.\par}
\only<3->{Groups: $6\div2!=3$. Both lists of a pair name the same members.\par}\end{overlayarea}
\end{frame}

% Frame 64. Teaching note: Ask what the divisor counts before advancing. No labels appear on the real matching copies.
\begin{frame}[t]{Oral check 3: matching copies}
\Q{Two matching circles and one triangle: how many visible rows?}
\DiagramStart
\Obj{C}{5.0}{2.6}{}\Obj{C}{6.1}{2.6}{}\Obj{T}{7.2}{2.6}{}
\DiagramEnd
\begin{overlayarea}{\linewidth}{1.15cm}\only<2>{$3!\div2!=3$. Each visible row has two circle-label versions.\par}
\only<3->{Every supplied copy appears once. The triangle has three possible positions.\par}\end{overlayarea}
\end{frame}

\end{document}
